\section[第二章\q 概念解释]{第二章\\ 概念解释}

\subsection{冰山原理}

{\kaishu
我们永远不可能认识事物的全貌，那么我们由于对事物不完全认识而形成的信念，也只不过是断章取义罢了。

\hfill ——金融帝国
}

我们先来思考一个问题，能够影响市场的因素是有限的还是无限的？按照混沌理论的说法，世界上的任何事物都能通过数次放大，从而影响其他“不相关”的事物。这就是“城门失火，殃及池鱼”的道理。如果说南美的一只蝴蝶扇动翅膀，可以导致北美形成一场飓风，那么一只蝴蝶煽动翅膀，同样可以导致农产品期货的剧烈震荡。也许宇宙空间的某个天体（太阳活动）发生一些小小的现象，就能很快地影响我们的金融市场。由此看来，能够影响市场的因素应该是无限的。进而我们可以得知，任何人对市场的分析、预测、理解都是片面的，都是断章取义的。毕竟我们谁也不可能考虑到无限多个因素。

这就涉及一个数学问题：自然数和偶数哪个数量更多？如果在有限的数字面前，那么自然数的数量无疑是偶数的2倍。但可惜数字是无限的，任何自然数乘2都能得到一个偶数。那么最终的结论就是：自然数和偶数是一样多的！

这说明什么？也许读者们会觉得这和交易没有任何关系。错，世界上的任何事情与交易都是有关的。在数学上考虑有限系统与无限系统的思维方式是不同的！我是说，交易者通常都是用有限的思维方式来分析无限的市场。如果能够影响市场的因素是无限的话，那么对于交易者而言，无论你是考虑到1件可能影响市场的因素，还是考虑到1万件可能影响市场的因素，在无限的市场面前，你考虑的都是一样多的！

我们都玩过“找茬”的游戏，就是找出两幅图画中有几处不同的地方。当我们找出一处不同的地方以后，图画中不同的地方就会少一处。直到我们找出所有的不同之处，两幅图画间就不再有不同之处了。这就是有限系统的思维方式，当你找出的因素越多，你就越接近真相。

一位做软件的朋友给我讲过一个软件界的理论，就是软件可能存在的缺陷数量等于已经发现的缺陷数量的N倍。这说明我们无法用排除软件缺陷的方法，来降低软件中可能隐含的缺陷数量。相反，你找到越多的缺陷，就越说明软件中会有更多的缺陷你没有找到。

既然无论软件设计者如何寻找软件中的缺陷，都无助于他得到完美软件的话，那么交易者无论如何综合分析影响市场的因素，都无助于他更接近于市场的真相。

在我看来，交易者所喜欢的“综合分析”，不过是因为本能排斥不确定性，而迎合本能来制造事后合理的幻觉。这方面内容读者朋友可以阅读一下“树形原理”的部分，你会发现事后的合理不等于事先的确定性。你考虑的因素越多，你就越容易事后解释行情。但同样你考虑的因素越多，你在事先预测时出现的矛盾也就越多，甚至你考虑到错误因素的可能性也将越大。

那么我们该如何应对市场呢？一座漂浮在水面上的冰山，我们无法看到冰山的全貌。我们能够看到的冰山越大，那么同时也说明水下我们看不到的冰山部分更大。我们看到的永远只能是冰山的一角。无论是冰山的最高点偏离了水面多高，最低点偏离了水面多低，同时也无论是更多的冰块融化成水，还是更多的水凝结为冰。请注意，水面的高度是不变的！

对于交易而言这至少说明：无论你找到多少可能影响市场的因素，而得出与市场现实（价格）有多么大背离的结论，都只能说明，你忽略了更多的相反因素能够与之平衡，并且使得分析结论与市场的真实价格相等。曾经在股市上，我习惯于买入不停上涨，但业绩很差的股票。当我已知的因素与现实的走势出现很大偏离的时候，我更乐意相信是我未知的因素在起作用。通常当我还不知道这只绩差股为什么连续不断上涨的原因之前，我相信立刻买入还是安全的。也就是说，通常我们未知的因素比已知的因素，更容易对未来价格的走势产生作用。

这就是我曾经在网上多次发表的“A+X=B”公式的原理。无论我们对市场的分析有多么详细，我们找到的原因仍旧是有限的，并且所有能够影响市场的原因还是可以分为已知原因和未知原因。如果承认结果，也就是真实的价格是由众多原因产生的话，那么就可以得出“已知原因A+未知原因X=结果B”的公式。通常对于未来走势起更大作用的，恰恰是未知原因。

当然前面已经阐述了，能够左右市场价格的原因是无限的。即使因果关系成立，那么无限的原因，也并不能得出一个唯一确定的结果。本公式所说的全部原因并不是能够影响市场未来的全部原因，而是现在市场全部参与者所能找到、考虑到的全部原因。由此可以看出，能够预测未来的公式并不存在，或者说这个公式即使存在也是具有不确定性的。至于未知因素更容易对未来走势产生作用，也是站在我们这些市场弱势群体的角度阐述的。

\subsection{树形原理}

{\kaishu
当你将一个球抛向墙壁时，通常球会被弹回来。但如果你能尝试足够次数的话，小球总会有一次能够穿透墙壁而出现在墙壁的另一边（因为布朗函数很小，所以你需要尝试的次数将是一个天文数字）。

\hfill ——量子理论的预言
}

我们从特定的一片树叶向下寻找，就必然能够找到唯一的树干；但如果我们从一棵特定的树干向上寻找，就绝对不可能找到唯一的树叶。如果你某晚回家输得连手表都不剩，那么你的太太肯定知道你是去赌博了；但如果你某天出去赌博，那么出门时你的太太无论如何也不可能知道你今天到底会输还是会赢。

为了更形象地说明树形原理要表达的含义，我们还是举一个更具体的例子吧。我曾经做过十年的交易，在这十年中我经历了无数的失败与成功。我的每一次经历，或多或少都会改变我的一些思维方式。比如，600065的惨败让我明白了止损的重要性；600788的惨败让我明白了走势说明一切；600745让我彻底放弃了对确定性的追求……也许这些感悟并无高深可言，但如果没有亲身经历的话，我想任何交易者都很难真正理解其中的含义。请注意，这些都是过去的事情，都是一些已经固定下来的事情，并且我现在的信念无不留下了这些经历的烙印。那么十年、二十年以后我在交易上最终的结局会是成功还是失败呢？这是未来的事情，并且是不确定的事情，而最终的成败必然很大程度上取决于我的信念。

我相信，如果未来我最终在交易上成功的话，当我回忆这些经历的时候，肯定会感慨曾经经历的每一次失败都只不过是在为我最终的成功做准备。如果缺少了曾经某一次重大的失败，那么我就会缺少了一个重要的感悟，最终我就有可能不会走向成功。2006年我在期货市场上用了两个月的时间，把17万变成了60万。那时我回忆曾经的一幕一幕，真的有上面所说的那种感受。甚至我当时说下一句豪言壮语：“至今对我一生影响最大的四笔交易中，有三笔亏损了超过50％，一笔获利了250％。如果让我舍掉其中一笔交易的话，那么我会舍掉这次获利250％的交易。”因为我知道，没有曾经的失败，我也不会有今天的感悟。

但我更相信，如果未来我最终在交易上失败的话，那么当我回忆曾经这些经历的时候，我会感叹道：我一生经历的一切，都只不过是在为了我最终的失败做准备。无论如何我的信念都要为最终的失败负责。失败虽然并不能等于我的信念是错的，但至少能够说明市场并没有符合我的信念。如果我相信的是其他的信念，那么也许就不会出现失败的结局。对于巴菲特的成功而言，我想一个不可或缺的因素就是市场符合了巴菲特的信念。如果巴菲特成长阶段的历史时空是在中国的1997年至2001年的话，估计世界上就不可能有什么巴菲特的存在了。

明白我说的意思了吗？历史是确定的，而未来是不确定的。同时不论未来最终是什么结果，我们都能用已经确定下来的历史进行合理的解释。但是，无论如何我们都不可能用已经固定下来的历史来预测未来的结局。所谓树形原理，说的就是原因与结果之间的非线性关系！简单地说，事后合理不等于事先可知。或者也可以理解为，“同样的原因可能产生不同的结果”。无论你使用什么样的分析方法，同样的原因都可能产生不同的结果。在我看来，交易者是否能够长期获利，很大程度上取决于如何应对这种不确定性。

我只想说，原因与结果之间是具有非线性关系的，而确定性的结果根本就不存在。如果非要说明是什么产生了非线性关系的话，那么恐怕就不是我力所能及的事情了。从上个世纪爱因斯坦与量子论众多创始人展开“猫的悖论”的争论，到最近刚刚兴起的混沌理论，都在反复对“未来是否是确定的”这个哲学话题进行探讨。但似乎有一点是能够达成共识的，那就是对于观察者而言未来是不可知的。具体这种不可知性是因为观察者无法考虑到无限的原因；还是因果关系从本质上就具有非线性的特征；或者是因果关系根本就不存在……这些似乎不是我们这些交易者应该考虑的问题。好在无论事实是什么，对于我们这些交易者而言都是完全等效的。抛开原理阐述现象的话，结论就是：无论你考虑到多少原因来推导结果，结果都不可能是单一的。但如果在结果确定的情况下，历史会给你留下足够多的原因供你解释这个结果。所以我形象地把这种现象称之为“树形原理”。

回到市场上来，我们可以经常发现这样的现象：交易者在回顾历史行情时，无论是利用基本分析、还是利用技术分析来对市场进行解释，总是为了维护本能需要的确定性，而通过因果非线性关系来制造一种幻觉。《通向金融王国的自由之路》一书中将这种现象称为“事后合理偏向”。从某种意义上来讲，无论是基本分析还是技术分析，能够对单次交易负责的方法都是基于幻觉。关于抛弃确定性对于交易的重要性，我们在稍后的章节中详细探讨。

\subsection{巨变效应}

{\kaishu
当一切都已经顺理成章，似乎永远都不会改变的时候，往往巨变就会随之而来。

\hfill ——金融帝国

很多人是靠偶然事件走向破灭，而我是靠偶然事件使资金上一个台阶。

\hfill ——大意来源于某国外书籍
}

在我们熟悉的市场上，经常会发现这样的现象。某一特定时点的前后，市场的特征变得截然不同。这就是我们通常所说的转折点，当然这个转折点可能是上涨与下跌的转折点；也可能是趋势与整理的转折点；还可能是任何信念的转折点。在1997年当业绩的信念已经深入人心的时候，当市场上流传着，买入长虹、发展只会赔时间而不会赔钱的时候。突然间绩优股开始陷入了长达四年的灾难，而曾经没有人要的“死亡板块”却走出了数倍、数十倍的上涨。由此可见，任何相互对立的概念都会交替被市场所接受，而往往这种转变是非常迅速的。

我经常会遇到股民指着股票历史走势图中某一箱形整理的阶段，告诉我在这个阶段应该暂时放弃趋势思路，而采取高抛低吸的策略。我通常会对他们说，整理阶段之前肯定是趋势行情，我们以上涨为例：一次上冲就见顶为尖顶；两次上冲见顶为双顶（M头），三次上冲见顶为三重顶或者头肩顶。这说明整理行情在走出之前，你不知道它是整理行情。而当整理行情真正走出来以后，通常就意味着整理行情可能突然结束。也就是说，更多的交易者都是根据短期经验来寻找规律，并且迅速死在突然的巨变中。从江恩的4次法则我们可以得知，通常第四次的上冲会冲出箱体。按照我个人长期使用的趋势跟踪系统来看，只要市场整理的幅度与我的趋势跟踪尺度正好吻合的话，那么我的系统就会出现使我亏损的假信号，而通常来说超过四次的假信号非常少见。当以逆势的敢死队思维进行高抛低吸的交易者已经出现获利效应的时候，你跟他们学只能是为他们来买单。通常我会使用相反的方法来应对这种情况。这方面的内容，我们在后面的“学习效应”中再进行探讨。这里我要说的是，当一种现象已经成为一种定式的时候，往往巨变就会随之而来。

{\kaishu
上个世纪初“历史上著名的拐错弯”把世界从似乎永恒的平静中引向了可怕的第一次世界大战；上个世纪80年代，苏联已经成立了75年，在世界舞台上似乎成为一个永久的国度，然而，让人愕然的是，短短几年后苏联就解体了；日本是个地震多发的国家，但在神户地区几个世纪以来却是一直非常平静。上个世纪末的一天，一场震惊世界的神户大地震使得巴林银行走向破产。（本段摘自《临界》一书）
}

当然，事后分析起来，这些偶然事件正巧碰上了合适的土壤，所以才演化成一场巨变。按照《临界》中的说法：大的事件并不一定需要大的原因。那么为什么这种巨变是迅速的，而不是缓慢的呢？我们以地震为例，地震的成因非常简单，就是由地壳运动而产生的。地壳的运动是非常缓慢的，而地震却能在瞬间释放出巨大的能量。这是因为岩石之间具有“摩擦力”，地壳运动产生的能量被这种摩擦力封存起来。当这种摩擦力与封存的能量处于一种临界状态时，非常小的偶然事件就能产生巨大的能量释放。

即使是在我们的生活中，我们也不难发现：原因通常是渐进的，而结果却是跳跃的，也许这就说明原因和结果之间确实存在着“摩擦力”。记得小时候我得过一次肺炎，最后好像是吃了一种在当时很新的抗生素才治好的。当时我老妈跟医生说：“还是西药的效果更好，以前吃了这么多中药都没有任何改观。”医生说：“我们吃饭都是吃了最后一口才吃饱的，难道我们之前吃的那么多饭都是没有用的吗？”在《短线交易大师：工具和策略》中我也看到了这样一种说法：交易者的成长就像一颗种子成长为花朵一样，起初相当长一段时间的成长都是在我们无法观察到的地下进行的。“不要不满现在，成就往往就在一瞬间”，说的也是这个道理。努力是渐进的，而成绩却是跳跃的。

上面说到的“摩擦力”本质上究竟是个什么东西呢？不地震是一种常态，而地震是一种非常态，那么摩擦力就是一种维持常态的力量。或者说常态是一种普遍的规律，而“摩擦力”又表现为现有规律的惯性。“摩擦力”的本质似乎还能解释为封存效应，我们来看看政府调控经济的效果吧。1997年中国经济成功实现软着陆以后，就进入了长期的通货紧缩状态。政府采用了包括激进财政政策、降息等一系列的刺激经济政策。当时这些政策到底起到了什么样的效果呢？我个人感觉几乎没有任何效果。2004年的时候，中国的物价才开始见底回升。现在想想看，当时房地产的价格真是便宜得不可思议。那么数年刺激经济的政策真的没有效果吗？不是这样的，刺激经济的政策虽然基本上不能扭转通货紧缩的现实，但却为未来的通货膨胀埋下了伏笔。上海的房地产一夜之间陡然而升，并且迅速蔓延至全国。短短两三年间，我身边的二手房价格上涨了五倍！调控经济就像是煲鸡汤，淡了加点儿盐，咸了兑点儿水。但遗憾的是，现实世界中存在着封存效应，盐并不能很快溶解，水也不能很快稀释。从微观上看，虽然政府在不断加盐，但汤反而越来越淡。突然，汤又变得咸得让人受不了，政府只好再不断兑水，但兑水后也改变不了汤越来越咸的现实……也许经济发展的实质就是在管理层不断加盐、不断兑水中，经济总量变得越来越大。

回到我们的金融市场上，巨变效应也是同样被经常演绎。当市场的常态已经被交易者所深信不疑时，通常巨变就开始了。市场的本质也许就是两种相反力量的正反馈与负反馈。上涨意味着继续上涨，就是正反馈；而上涨意味着下跌，就是负反馈，同时负反馈具有封存效应！上涨的本身就在不停地聚集下跌的力量，同时上涨又让交易者对下跌的理由视而不见。当市场最终处于临界状态时，一个小事件就能使聚集起来的下跌因素瞬间爆发出来，从而产生巨变现象。外国某交易大师曾经说过：大多数的基金经理在大多数的时间内都是正确的，而在少数偶然事件面前输得一败涂地。如果你真的了解市场的本质，你会发现这种偶然事件，只不过是通常很少出现的必然事件！我深信：很多人是靠偶然事件走向破灭，而我是靠偶然事件使资金上一个台阶。

\subsection{恐龙效应}

{\kaishu
交易者戴上光环与走向陨落，从本质上来说其实是同源的。

\hfill ——金融帝国
}

7800万年以前地球的自然环境是那么温暖和潮湿。此时爬行动物非常迅速地适应了当时的自然环境，将自己快速发展壮大起来。由于对环境的适应，而抢占了越来越多的生存资源，最终成为当时地球的主人。突然地球的自然环境开始变得寒冷与干燥，曾经带上光环的恐龙迅速地灭绝了，而从来就没有展露过锋芒的昆虫却活了下来。

市场和我们生活的自然界有很多相似之处，在不同的时间段内表现出截然不同的市场特征。以股市为例，1995年以前市场热点是小盘、绩优、含权；1995年市场热点是异地、大盘、低价；1996年到1997年上半年市场热点为绩优、大盘的行业龙头；1997年下半年到2001年市场热点为绩差、小盘、曾经被誉为“死亡板块”的题材股票；2001年至2008年的市场热点又回到以业绩为主线的价值投资。

大盘与小盘、绩优与绩差、价值与题材等等，这些互相矛盾的概念都曾在股市某个特定的阶段称霸一时。1998年至2001年市场以不炒市盈率而炒“市梦率”为主流理念，而之前和之后市场的主流理念都在宣称“业绩是市场永恒的主题”。从基本面来划分的话，还可以把特点的概念进行更细微的划分。但这些更微观的概念，似乎寿命会更加短暂。

从技术面来看，市场在不同的时间段内也表现出截然不同的市场特征。有时价格在一个非常窄幅的空间内不断整理，有时价格又不加停顿地勇往直前；有时市场牛市冲天，有时市场又熊气弥漫；有时市场运行得非常缓慢，总是一步三回头，有时市场的趋势特征又强大的超过所有人的想象力。

根据“巨变效应”我们可以发现，这种截然不同的转变往往发生在非常短的时间内，而并不是缓慢渐进地改变。2001年6月之前市场不温不火地横向整理，而之后市场却迅速走出了夸张的下跌行情；2005年7月以前市场熊气弥漫，而之后似乎突然就变成了势不可挡的牛市。这种在两个极端之间跳动的现象，在市场上屡见不鲜。

做交易不外乎就是靠对规律的寻找，找到某种概率优势。而规律我们又可以分为微观规律、宏观规律。越微观的规律就越能让交易者创造出辉煌的战绩，但同时越微观的规律能够保持的有效时间就越短。同样，越宏观的规律就越能够保持更长的有效期，但同时越宏观的规律就越难在某个特定的时间创造出辉煌的战绩。

曾经在股市上，熊市中我通常不是做得最好的，做得最好的是死空头；牛市中我通常也不是做得最好的，做得最好的死多头；但有意思的是，一轮牛熊循环过后，能够超越我战绩的人总是寥寥无几。所谓死空头就是把熊市中任何上涨都视为反弹的人，死多头就是把牛市中任何下跌都视为回调的人。而对于我来说，无论牛市还是熊市，也无论我是看涨还是看跌，只要是能够超过一定幅度的反向运动，我都会视为趋势的改变。如果趋势没有改变，我会在其后更不利的价位重新找回自己的位置。很显然有成本的人，永远没有无成本的人利润大。

也许“恐龙效应”真的能够很好地解释，为什么很多交易者都是在戴上光环后快速地陨落。能够戴上光环的交易者，必然更多地利用、依赖了市场中的微观规律。同时他们的辉煌也不可能过多地在交易中考虑到成本问题。一个交易者越能适应当时的市场环境，似乎就越不可能适应变化后的市场环境。这时通过过度适应市场而戴上光环的同时，也就为其后的陨落埋下了伏笔。当市场的特征出现巨变的时候，这种过度适应市场的交易者就会被市场迅速扫地出门。

在我们的地球上生活着脱颖而出不久的哺乳动物，也生存着远为古老的昆虫，但这之间的恐龙却灭绝了。有意思的是，在我们能够看到的交易书籍中，也存在着别无二致的现象。我们能看到最近几年才风靡起来的最新交易技术，也能看到上个世纪初流传下来的交易理念，但这之间的几十年似乎真的是一种真空。记得1997年时，市场流传着一种“小波浪理论”的分析方法，据说当时抄底逃顶屡试不爽，但后来这种方法就没有任何用武之地了。也许越古老、越简单的交易理念才能越接近交易的本质，而那种精确的交易技术似乎总是会迅速被人遗忘。

按照这种理论推导下去，我们不难得出这样的结论：好的交易者应该是那种从来没有展露过头角的人。这和我们的常识似乎不太相同，交易比赛的获胜者在展示他们辉煌战绩的同时，总是能够吸引更多交易者的眼球。在我看来，好的交易员应该是默默无闻的，他们没有任何可以炫耀的辉煌。而这些默默无闻者，却一遍一遍地看着其他高手从戴上光环到迅速陨落的表演。

\subsection{钟摆效应}

{\kaishu
市场并不存在为单次交易负责的单向真理，市场上的真理在概念上都是成对出现的。

\hfill ——金融帝国
}

我小时候在学校内看一本关于股市的书，有一位同学觉得挺好奇，也非要看一看。由于他并非“身在此山中”，所以也就没有股民身上固有的那种追求幻觉（完美的确定性）的倾向，从而说出了一句非常经典的话：“我发现书中的内容几乎都是矛盾的”。顺便说一句，那不是一本关于一整套交易理念的书籍，而是一本东拼西凑起来的有关交易技术的书籍。对于现在的我而言，这样的书根本就不值得一看。

没错，我这位根本没有做过交易的同学，真是说出了关于市场的真谛。市场当中根本就不存在什么能够为一次交易负责的单向真理，而是由众多对立矛盾的概念组合而成的。我们这些交易者都是在这样的至理名言中成长起来的：要顺势而为，但不要追涨杀跌；要高抛低吸，但不要逆势交易；交易者不能够贪婪、不能够恐惧……我就奇怪了，你不贪婪，干嘛还要来到股市？

市场至少从概念上来看是对称的，贪婪与恐惧、趋势与整理、敏感与迟钝、业绩与题材、大盘与小盘等等都是“成对”出现的。当你能够找到一种交易的方法，那么我肯定还存在一种与你相反的交易方法。或者说，任何方法在不同的历史时段内都会有“占便宜”和“吃亏”的时候。有人认为，业绩是股市永恒的主线，但这种思路在1997年至2001年的4年间将会陷入泥潭。我深信，市场就像一块磁铁，有S极就有N极，而磁单极并不存在。

当然上面的论述完全是站在宏观的角度上，而在微观上市场确实能够在某一特定的时段内重复着某一极端的单向规律。根据“巨变效应”中阐述的内容，我们可以想象得到：市场在宏观上应该是在两个单向的极端间摆来摆去，这就是我所说的钟摆效应。

市场的钟摆效应确实存在着钟摆的某些特征。我们经常听到这样的说法：横的时间有多长，立起来就有多高。这就有点类似于江恩理论的“时间换空间”，市场用整理的时间换取了未来大幅波动的空间。同时也有“空间换空间”的说法，市场用大幅上涨的空间换取了未来大幅下跌的空间。在江恩看来，时间和空间从本质上来说是一回事。从交易者的信念中，我们也能够看出这种转换现象。当1997年业绩信念深入人心后，市场进入了长达4年的题材市；而当2001年题材为王的信念深入人心后，市场至今都在围绕着业绩的主线。这种“捧得越高，摔得越狠”的现象，与真正钟摆左边摆动得越高，右面也就摆动得越高的规律不谋而合。

如果说钟摆的摆动是绝对对称的话，那么市场是否也是绝对对称的呢？假设市场是绝对对称的话，那么任何一致性的交易方法在交易足够多的次数以后，最终的战绩都会趋向于零。这时任何交易方法都不可能在宏观上从市场上获利，而不采用基于从概率引申出的一致性交易方法的交易者又和赌博无异。这时的市场就是绝对的效率市场！

我不否认市场在很大程度上是对称的，是具有“效率”的。但万幸，市场并非完全对称。比如，经济会随着时间而不断增长，这就是一个不对称性。由此就能够产生一种“随便买、随时买、不要卖”的长期指数化投资策略。市场是人性的一面镜子，如果人性具有不对称性的话，那么市场也就具有不对称性。本书的全部交易理念都是基于人性的不对称性，使得市场具有某种概率上的不对称性。利用这种概率上微弱的不对称性，建立一种具有概率优势的一致性交易策略，从而实现长期必然的获利。这些我们会在后面详细地探讨。这里我们需要了解的就是：虽然市场并非是完全对称的，但对于一种方法而言，前人的利润必然会由后人来买单（当然这种买单未必和之前的利润绝对相等）。

这种不完全对称是存在无疑的，但同时又是非常微小的。既然股民的战绩具有“7赔、2平、1赚”的规律，那么市场就绝对不可能是完全对称的。有意思的是，股民并非次次出错，所以这种不对称性也就是略微的。在市场上长期必然获利的方法，完全来自于大多数股民的集体误区。这样一来，我们如果能够研究明白大多数股民为什么亏损，那么就能轻松找到获利的办法。

\subsection{高低错觉}

{\kaishu
强加给观众一个参照物，那么观众就会得出一个你所需要的高低评价。

\hfill ——金融帝国
}

在我看来，更多的交易者都不是趋势思维的信徒。他们不甘心当“傻子”，而在任何一笔买卖当中都要考虑所谓高低的问题。这本身就是一个误区，高低根本就不应该作为买卖的理由！对于我这样的趋势信徒而言：只要上涨还在继续，那么多高都不嫌高；只要下跌还在继续，那么多低都不嫌低。可是市场上总是流传这样一种说法，交易获利的根本在于低买高卖。也许这种说法并没有任何错误，只有未来卖出的点位比买入时更高，交易者才能够获利。请注意，这时的“低”是以未来作为参照物，而绝对不是以过去作为参照物！遗憾的是，更多的交易者都喜欢同历史比较而得出高低的结论。在他们看来，只有买的比历史低，才算是聪明。如果今天的价位比昨天高，那么今天买入的人就是傻子。

也许市场除了有趋势特征以外，确实还存在整理特征。这就是所谓的负反馈，涨多了就要跌，跌多了就要涨。那好，我们就同历史比较，来迁就整理思维的信徒吧。这时有一个不容忽视的问题：既然所谓高低是一个相对的概念，那么我们该如何选择参照物呢？

大盘从998点涨至6124点，然后又跌至4779点。那么这个4779点到底是高还是低呢？如果我们画一条水平线，就能够找到2007年8月13日的上证指数也是4779点。那么同样的4779点到底又是高还是低呢？我相信，大多数股民对于第一个4779点都会以6124点为参照物得出低的结论，而对第二个4779点又都会以998点为参照物得出高的结论。这时同样的4779点在不同的参照物面前，得出了截然不同的高低评价。由此看来，即使是整理思维的交易者，他们对于自己极其在乎的高低问题，所作出的评价也是不客观的！

股民通常都喜欢卖高买低，这时对于个股而言就会出现一个有意思的现象。股票越涨市场介入程度就越低，而股票越跌市场介入程度就越高。所谓市场介入程度，是指散户的参与程度。道理很简单，一支反复上涨的股票，曾经的持有者会因为价格高而卖出，同时持币者也不会买入找不到高点作为参照物的股票。我有点理解大多数散户了，自己当了傻子不要紧，只要之前有更傻的人就行。事实上，任何一个点位买入股票，都有可能被套牢。但是买入上涨的股票就有可能成为最傻的人，而买入正在下跌的股票就会认为之前的买入者肯定比自己傻。这就是人性丑陋的一面，就像老话说的“死了也要拉一个垫背的”。当600788跌至0.2元时，14元的买入者要比15元的买入者好受一点吧。在2001年的散户厅内，我发现了这样一种可悲的现象。一群炒股票的中年妇女，互相打听对方套牢的股票，然后她们就去抄同伴的底。无论未来的涨跌，他都能占到同伴的便宜。后来的事情就更加可笑了，这些人在绝对公平、自由的资本市场上，居然开始勾心斗角、甚至关系紧张。

在那个庄股横行的年代，庄家反复利用散户这种心理弱点以及对高点的判断错觉。一只股票从5元涨到10元时，散户会越来越少，这时筹码流向庄家和趋势跟踪者手中。请注意，更多的散户不是趋势思维的信徒，所以无论股票怎么涨，敢于追高买入者都是少数，这样一来庄家无论如何都出不来。当股票一路上涨时，能够抛售的散户越来越少，所以越涨就越不费力气。如果庄家准备在20元实现成功撤庄的话，那么通常会拉高到30元、甚至40元。道理同样很简单，没有“打折效应”，散户是不会当傻子的！坐庄并不复杂，从5元到10元完成吸纳，然后空拉到40元，之后的下跌就会不断地套牢抄底盘！当然随着股票的上涨，基本面都会或多或少地变好，最低限度也会流传出来一些概念。可爱的散户朋友们，就会用折扣价格买入富有丰富概念的“优质股票”。按照常理来说，基本面通常是跟着股价走的。我们可以把基本面理解为60日均线，也就是说基本面的最高点通常落后于股价的最高点。当股价与基本面（60日线）相交以后，股民就会有一种超值的错觉！如果股民习惯于将过去三个月的价格当做参照物的话，那么60日均线就成为了股民对于高低判断的形象表述！

这不由得让我想起一种销售策略。数年前，我所在的手机公司推出一款具有“超长待机”功能的手机。那时这种功能并不普及，所以消费者也很难说得清其到底价值几何。如果这款手机低价销售的话，那么消费者就会认为超长待机并不值钱（其实也是不值钱）。公司的销售策略就是：花费重金宣传，同时高价销售。我们这些业务员都感到销量不大，建议公司降价销售。公司给我们的答复是：现在并非让你们“冲量”，而只要让消费者知道这款手机即可。只要我们把高价位挺住两个月，那么未来的每一次降价都能使销量上一个台阶。确实高明！消费者如果没有参照物的话，很难对于贵贱作一个“合理”的判断。就像股民也同样并不了解股票应该价值几何，尤其是不明白一种题材、一种概念价值几何，没有参照物同样也很难判断的道理是相同的！

这不由得又让我想起一个有趣的故事：佛祖某日下山看到一家商店供奉着一尊精美的佛像。佛祖非常喜欢，决定将其买下来。老板看出佛祖的喜欢之情，所以将并不值钱的佛像开价5000元！佛祖扬长而去，回寺与众僧讲述此事。众僧都谴责利欲熏心的商人，并且决定教训他一下。但佛祖说：这倒不必，他有办法用1000元将佛像买下，并让众僧听他的安排。每隔一日让一位僧人以比昨日低500元的价位去商店内购买。第二天僧人出价4500点，店老板不卖；第三天另一僧人出价4000，店老板同样不卖……直到数日后，一位僧人出价500元，店老板肯定还不会出售。次日，佛祖再度下山，轻松地以1000元的价格将佛像买下。这其中的奥秘并不复杂，读者朋友们可以想象一下店老板每日的心情变化。其实经历过熊市的读者，一定会深有感触。

在我看来，趋势与整理并不矛盾。事实上，每个交易者都在运用趋势和整理这两种思维，只不过他们的尺度并非相同。关键是，让整理的尺度大于趋势的尺度一个数量级！中线追涨杀跌，那么长线就是高抛低吸。遗憾的是，更多的交易者乐意于中线高抛低吸，长线追涨杀跌。如果股票从10元涨到20元的话，那么12元买入股票、右半侧跌至18元卖出，难道就不是高抛低吸吗？事实上，更多的股民乐意于12元高抛，而18元低吸。

\subsection{对错悖论}

{\kaishu
对错的评价不应该是基于结果，而应该是基于原因。

\hfill ——金融帝国

如果一位交易者学不会如何正确区分“对与错”，那么他将永远原地转圈而毫无进展。

\hfill ——金融帝国

在交易之路上，我们走得再慢都没有关系，但是我们必须确认自己是在向前走。

\hfill ——金融帝国
}

某日一位赌徒赢得了比一年工资还多的钞票，回家后太太像迎接功臣一样准备好了丰盛的酒菜。不久这位赌徒又输掉了更多的金钱，这次他的太太说什么也不让他进门了。赌博是原因，输赢是结果，是否赌博是可以确定的，而输赢是无法确定的。那么交易者对于交易方法而言，是否也像这位赌徒的太太一样，用结果来决定一切呢？

有时我在想：我们这些交易者的喜怒哀乐、甚至是对错的评价，都是完全建立在我们根本不能左右的结果上，这难道不可笑吗？在大多数的交易者看来，能赚到钱的方法就是好方法，而亏钱的方法就是坏方法。如果真的按照这样的逻辑来分析对错的话，那么任何方法都是好方法，同样任何方法也都是坏方法。

这种对错观的直接产物就是我们后面要谈的“学习效应”，这样一来我们就陷入了永无止境的“狗咬尾巴”的怪圈而无法自拔。事实上，交易者达到成熟的起点与终点之间的距离并不遥远，但通常交易者都是在原地转圈而停滞不前。在我看来，交易者能否取得成功，从某种意义上来说就是取决于坚持“错误”与改正错误的能力。其中第一个“错误”代表自身交易策略带来的合理亏损，而第二个错误代表交易方向错误而出现的亏损。这就是庄子所说的“外化而内不化”。这方面的内容我们稍后再详细探讨，这里我希望读者朋友们明白的是：对于对错的评价应该是基于交易策略而不是基于结果。

对于任何交易者来说，都必须找到自己应该承担的东西。交易绝对不是没有成本的不劳而获，所以我们必须首先确定我们能够承担的成本是什么？然后，找到一种合理的策略来一致性地实现全部交易。这就是每一位交易者成长中找到的基石！当你知道什么样的行情是你必然亏损的行情时，那么你自然就找到了什么样的行情是你必然能够获利的行情。在交易之路上，我们走得再慢都没有关系，但是我们必须确认自己是在向前走。如果你不能接受“盈亏同源”的事实，那么你就不可能找到固定的基石，自然你的全部交易技术就只能是空中楼阁。

在“巨变效应”中，我们讲述了“原因是渐进的、结果是跳跃的”的道理。一根单独的稻草无法压倒骆驼，之前的稻草并非毫无价值。没有压倒骆驼的稻草就是你的亏损；而最终压倒骆驼的稻草就是你的利润。这就是“盈亏同源”的道理！当你规避了亏损，自然也就规避了利润，明白吗？最后一根压倒骆驼的稻草并非是特殊的！那么，让你盈利的交易也并非是由于特殊的方法而产生的！请记住，盈亏是不可分割的！

\subsection{学习效应}

{\kaishu
市场是一位狡诈的老师，而交易者是一位愚钝的学生。

\hfill ——金融帝国
}

向成功者学习是人类的一项重要本能。就连这一点，在市场面前都能成为不小的问题。有时我在想，到底是人性的错，还是市场的错？也许市场本来就是人性的天敌。无论人性如何，市场都会成为人性的敌人。

如果你没有由固定基石形成的一致性交易策略，那么你自然会受到身边交易者的影响。在一个大户室中，你左边坐着追涨杀跌的A，右边坐着高抛低吸的B。如果市场走出持续的整理行情，那么B就会获利丰厚；而A就会不断地左右挨打。由于你没有跳出“结果决定论”的怪圈，那么你就会认为B做的比A好。自然你就更倾向于向B学习……当你也决定采用B的方法来不断获利时，突然间市场走出了强劲的趋势行情。很快B送回了大部分的利润，而你只好平白无故地产生亏损。这时，A终于不再垂头丧气，而是昂首挺胸地享受他的利润，以及众人羡慕的目光。开始你觉得这是偶然现象，但趋势却一点也没有完结的迹象。一段时间后，你终于改变了自己的看法，你觉得事实上A的方法更好！好，这时你又开始跟A学。但更有意思的是，市场突然又陷入了似乎没有尽头的整理行情中……

后来你决定不能再受外界的干扰，自己用趋势跟踪系统来应对行情。开始时，你用一个敏感的尺度来捕捉趋势。本来你做多，但股市偏偏开始回调，并且触及你的止损点。要命的是你刚一卖出，市场就立刻上涨。经过自己的总结与反思，你发现是你的止损点设置得太过敏感。过了相当长一段时间，迟钝一些的止损点又一次让你卖在地板价上，你觉得还应该再迟钝一些。又过了相当长一段时间，你发现虽然股市走出了趋势，但你“掐头去尾”后并没有赚到什么钱。如果你能够敏感一点的话，那么就能赚到更多的利润……

上面两个故事的主人公就是曾经的我！请注意，这只是趋势与整理、敏感与迟钝两个对立的极端。事实上，市场中还有更多矛盾的对立。如果你在每一个矛盾的对立两极间，都转两个圈，那么十几年就过去了！事实上，从1997年到2001年我还从绩优概念摇摆到题材概念。

整理和趋势之间是不完全对称的。趋势略微强于整理，这是由股民更乐意于“盈利兑现、亏损挂起”而决定的。至于敏感与迟钝之间就是完全对称的，如果你用不同的尺度操作足够长时间的话，那么战绩是趋同的。当然这需要忽略了交易成本，但交易成本对于中线以上的尺度而言，基本上可以忽略不计。当你在市场上待的时间越长，你就能够找到越多的“等效性原则”。如果你总是无私地付出了成本就离开，而绝对不享受利润的话，那么亏损已经就是注定的事情！

市场的机制就是制造微观定势，然后突然打破这种定势，将追随定势的交易者一网打尽。这有一点像狼来了的故事。通常来说，微观定势每重复一次，信徒就会成几何倍数增长！牛市这位老师就是在不断地教给学生：止损是愚蠢的、长期投资是明智的。当好学的学生学会这些以后，熊市这位老师就开始上场了，他不断地教给学生：止损是英明的、长期投资是胡闹的。而且市场这位老师又是尽职尽责的，学生学不会、他就不离去。真是“包教包会、学不会不收费”。大牛市中的洗盘其实有两个作用：一个是通过持有者的换筹来提高市场平均持有成本；另一个更为重要的作用就是告诉大众——下跌并不可怕！

这个世界上只有两种人最吃香：一种是信仰坚定的人，一种是根本就没有信仰的人。我对价格的涨跌就是绝对没有信仰的人，绝对的随势而动。但我对趋势思维又是信仰异常坚定的人，无论输了多少次，我都还敢继续下注。而如果你是既有信仰，但又不坚定的话。那么你就会随着市场这座大钟而左右摇摆，并且永远慢市场半拍。先来的人吃饭，后来的人买单，你将永远是没有吃饭，但却买单的人。

我见过很多炒股时间超过十年的人，大多数仍旧无法跳出“狗咬尾巴”的怪圈。也许学聪明简单，但坚持愚蠢就不是那么容易的事情了。有的时候，敢于同最近的成功者说不，真是一种非凡的能力。我曾经深深地感受过那种巨大的压力，那种与潮流对抗的压力。

人通常都是健忘的，而只是对最近半年左右发生的事保持深刻的记忆。市场的评论家们总是能够找到最近半年获利最多的交易方法，然后张大嘴巴告诉所有的人。守住自己的“道”真的是一件很难的事情，要知道我们每一个人都不是孤立存在的。

有时我也在想，为什么交易上的成功者如此少？首先，没有任何人会教给投资者抛弃确定性的成本观念。好，大家都想次次正确，那么分析师就不敢谈不确定性、不敢谈赌博。其次，市场上的一轮循环可能持续极长的时间。热点从业绩到概念，再从概念到业绩的循环真的需要长达尽十年的时间。即使你再有悟性，没有经历过一次循环的人也是不可能感受到市场的“钟摆效应”。最后，就是人性的问题了。没有人会看到别人赚钱不“眼热”；没有人愿意心甘情愿的吃亏；没有人愿意相信单笔交易是没有道理可言的。

\subsection{超市原理}

{\kaishu
在我看来，股民们大面积亏损的一个最为主要的原因就是正确率太高了！

\hfill ——金融帝国
}

超市那种最后一次性付账的销售模式更能激发消费者的消费欲望。如果是每拿一件商品就付账一次的话，那么估计销售额就会出现下降。简单地说，人对于“次数”要比“数量”更为敏感。如果是收入的话，那么10次1元的收入产生的满意程度大于1次10元的收入。同理，对于支出而言，10次1元的支出产生的痛苦程度大于1次10元的支出。

这一点在市场上表现的也极其明显。为什么股市上牛市中的成交量大于熊市呢？因为交易者在获利的情况下更乐意交易。这同时说明交易者并不是绝对客观的，而是存在相当大程度上的偏向性。也就是说，交易者更重视概率而不是赔率。

虽然市场上有“7赔、2平、1赚”的说法，但市场平均正确率却是高于50％的。我们虽然不可能精确计算市场的平均正确率到底是多少，但用一个简单的公式就能算出一个大概。用股市上牛市成交量的总和除以全部成交量的总和等于全部交易者的平均正确率。有一个奇怪的现象：我在身边的股民中，找到一个正确率低于50％的交易者是非常困难的！如果你不相信的话，也可以试着找找看。

在我看来，股民们大面积亏损的一个最为主要的原因就是正确率太高了！假设某股民在最短持有周期为两周的情况下刻意降低正确率，那么就有可能实现长期的盈利。这又不符合大众的常理，难道亏损的罪魁祸首真的在于高正确率吗？

正确率的本质到底说明什么？正确率的高低，同预测水平的好坏并非具有绝对的相关性；而与离市策略具有更大的相关性！理论上讲，即使在随机开仓的情况下，只要止损的幅度数倍于止赢的幅度，那么正确次数就能数倍于亏损次数。所谓高正确率只能说明利润是被众多的人（次）平分，而亏损是由少数人（次）承担。由此看来，股民的高正确率，只能说明“利润兑现、亏损挂起”的癖好在起作用。简单地说，对于一定数量的利润而言，多数交易者希望是由更多的盈利次数与更少的亏损次数组成；而我希望由更少的盈利次数与更多的亏损次数组成。

交易者更愿意锁定利润，而逃避亏损的本能，正好为我们找到了获利的武器。我再说一遍，所有具有再现性概率优势的交易方法，完全是来自于人性的不对称性使得市场产生的不对称性。像大众对待利润那样对待你的亏损，你就能长期盈利！

\subsection{复利悖论}

{\kaishu
无论你承认与否，交易者追求的很多东西是永远也无法真正得到的。

\hfill ——金融帝国
}

如果我们能够每年获得40％的收益率，那么7年后资产就增长了10倍以上，20年以后资产就增长了将近1000倍。如果半年能够获利40％的话，那么10年后资产就能增长将近1000倍，而20年就能增长1000000倍。

年收益40％相当于每月平均获利2.8％，这是够低的收益率了。但如果你能够将这种微小的月收益保持20年的话，那么10万元的初始资金也能够变成一个亿！而如果你能够每年将资金增长100％，那么10万元的初始资金20年后能够变成1000个亿。要知道每年100％的收益率，只相当于每月6％的收益！

在我看来，20年1000倍的投资收益理论上来讲是可行的，但无论如何20年100万倍的投资收益是绝不可能的。那么由此可以得知，一年超过1倍的收益如果能够稳定的话，那么就能产生绝对不可能出现的“市场怪物”。由此可知，一年超过1倍收益率的交易方法，必然是不稳定的！要知道一年1倍的收益率只相当于6％的月收益，那么每月能够产生超过6％收益率的交易方法同样是不稳定的。

衡量获利速度的快慢有三个指标：正确率、赔率、交易周期。从“超市效应”中我们可以得知，正确率同赔率是互损的。即使全部交易环节都采用随机操作，那么正确率与赔率的乘积应该等于50％。如果我们能够通过任何交易技术将期望收益略微提高，那么我们就可以具备期望收益优势。而这种期望收益优势在反复的操作中被不断地叠加，那么就能产生真实的钞票。如果这种操作周期足够短的话，那么就可能在短时间内创造出惊人的利润！而这种惊人的利润在长期的复利效应下，就能产生根本不可能出现的“市场怪物”。

如果站在市场怪物不可能存在的前提下，那么越短的交易周期，就必然会使期望收益降得越低。同时我们略微的期望收益优势，还会被交易成本所大量吞没。但这与事实并不相符，我们可以轻松找到短线快速获利的案例。那么结论就只能是获利速度与稳定性互损！而稳定性说明了什么呢？稳定性可以理解为交易结果与理论值的偏差，或者是期望收益未来可能变化的幅度。在我看来，稳定性是最为重要的一个参数。它是决定交易者生死存亡的问题，稳定性差意味着被扫地出门的可能性更大。即使期望收益是正值，但只要交易结果偏差足够大的话，我们照样可能血本无归。如果说，交易者第一目标为“不亡”，那么我们就必须追求稳定性！

提高稳定性就必须降低获利速度，也就是说在期望收益不变的情况下提高操作周期。如果是期货市场的话，那么上面的说法还忽略了一个问题，那就是仓位！其他条件相同的情况下，越重的仓位就能产生越大的获利速度。既然期望收益是每位交易者所追求的，那么交易者的陷阱就只剩下重仓与短线！

如果我对一位交易者说，我正在开发一种能够在未来20年后将10万元变成一个亿的交易系统。那么这位交易者通常会认为我的欲望太高了，甚至认为我是在做一个不可能成功的系统。有意思的是，也许这位交易者认为自己的目标很保守，只不过想通过日内操作来实现每月2.8％的收益。按照这种逻辑，他认为每月获利2.8％比20年获利1000倍更容易！要知道，只要这位交易者能够在市场上活20年，那么他贡献的手续费已经是一笔不小的财富了。如果再加上移滑差价等交易成本，那么每月获利2.8％所需要的技术水平要数百倍于我。

顺便说一句，请真的不要小看手续费！从某种意义上来讲，市场上唯一能够绝对确定的事情恐怕只有手续费了。我与一些在期货公司上班，但不做交易的内勤工作者聊天时，探听到一些显而易见的内幕。第一，凡是能够经常在期货市场上赚钱的交易者，往往对公司的手续费贡献都是很少的；第二，交易真正做得好的人，往往都不盯盘；第三，超短线的交易者的确存在短期获得暴利的可能性，但死是肯定的，区别在于死之前能够放出多大的光。

我与很多人讲我在交易事业上的“十年大计”时，他们都说我太贪了。但有意思的是，我的交易技术很粗犷，而更多的人都追求很多精确的交易技术。我保守地计算了一下，如果大多数人追求的交易技术能够实现的话，那么他将10万元变成一个亿最多用五年。我是说，更多的交易者对于交易技术的要求同对交易战绩的要求，简直就是一天一地！没错，对于交易技术而言，更多的人都觉得我不思进取；但对于交易目标而言，更多的人都觉得我异想天开。而在我看来，他们追求的交易技术，是永远都不可能达到的幻觉！

很多交易者对于交易技术的选择，都抱着一种“射箭”的思维。无论我们最终能够射中几环，我们都要瞄准“十环”。事实上，交易根本就不是这样。如果我想射中“五环”，那么我就应该选择能够达到“五环”的交易技术，并且追求更高的稳定性，而最大程度地保证我能够达到这个目标。相信我，交易是一门“舍”的艺术。当我放弃了射中“十环”的可能性时，我选择射中“五环”的技术就有了更大的稳定性。

相信我，对于交易技术而言只有“可行”或“不可行”两种状态。当你准备使用一种交易技术时，可以先计算一下成功后能够带来多大的利润。如果这种利润都让你绝对不可思议的话，那么就不要在这种交易技术上过多地浪费时间。当你知道利润是不可能实现、而收入又是巨大的时候，那么只可能是有更大的成本隐含在其中。如果你连成本是什么都不知道，那么使用这种交易技术无异于自杀！

\subsection{鸵鸟现象}

{\kaishu
我是一位知道自己是在赌博的交易者，而更多的人是自以为在做交易的赌徒！

\hfill ——金融帝国
}

据说，当鸵鸟被猎人捕杀时，它们不是逃跑，而是把头钻在沙中。当他看不到枪口时，枪口就真的不存在了！中国人有一个特别不好的习惯，那就是回避问题。不愿意接受的事情，我们可以不去正视它，回避就是最好的办法。

我曾经受到过两次虚假业绩的伤害，差一点就让我血本无归。当时我真的痛恨这个市场太不公平！这样一来，我就陷入指责与牢骚中不能自拔……让我走出泥潭的是一篇和股市没有丝毫关系的文章——《宽容的荷兰人》。如果一种现象你无法消除它，那么与其回避还不如正视它，从而承认它、接受它、研究它、合理地控制它。如果你连承认都不敢承认，那么你是无论如何也不能合理控制它的。于是，我就承认并且接受业绩可能虚假这个现实，从而研究如何控制这种风险对自己的伤害。慢慢地，这种理念甚至成为我其后交易策略形成的一条主线，那就是：从对抗到包容。

至此，在我的词典里已经没有“忌讳”二字。我们忌讳的通常都是真实存在的事物，只不过我们不愿意接受所以才忌讳。那么这里就说一下大多数股民们忌讳什么吧。

首先，是“怕亏情结”。大多数股民都不乐意面对亏损，很少有股民在交易之初就计划好亏损、更不用说把亏损设计为交易策略的一部分。无论你是否承认，只要是做交易就有可能亏损！“盈利兑现、亏损挂起”虽然能够暂时对亏损视而不见，但通常未来都要承受更大的代价。既然与亏损对抗不会有什么结果，那么我就想办法包容亏损，用利润来对冲亏损。反正怎么交易都要亏损，那么我与其被动地接受亏损，还不如主动地、有计划地、有步骤地去亏损。当我完成买入（开仓）交易以后，全部的工作就是对亏损的控制。盈利嘛，不用管它、随它去吧，它会自己照顾好自己。现在的我真的很轻松，回头看看那些忌讳亏损的人，他们活得真是很辛苦。一天我在股市里等人，看到一位股民十分钟向他的同伴问了八遍“你说过两天能不亏损吗？”这种亏不起的心态，最终的结果只能是上涨没他事，下跌他肯定跑不了。其实学会止损只是初级标准，然后还要不忌讳亏损。当你明白“盈亏同源”后，你就知道亏损本身就是最好的风险释放过程，这时你会爱上你的亏损。如果亏损是你永远也不能征服的敌人，那么为什么不将亏损包容、然后和他做朋友呢？我们可以不用同利润亲密接触，但对于亏损我们必须形影不离地对他照顾有加。我绝对不是开玩笑，交易的真谛就在于像常人对待利润那样对待你的亏损。

其次，是“怕隔夜情结”。在期货市场上，由于很多品种的定价中心在欧美，所以晚上外盘的大幅度波动就很可能对交易者的隔夜持仓造成很大的风险。这种隔夜风险是确实存在的，而且我至今也没有找到这种隔夜跳动任何概率上的偏向。这样一来，在期货公司的怂恿下，更多的交易者都把不持有隔夜单作为自己的交易规则。道理很简单，很多交易者都曾经受到过隔夜跳空的伤害，自然就不愿再涉及这块伤心之地，而是陷入日内和重仓的误区。在日内的有限时间内，交易者基本上不可能靠放长利润来实现大赔率。自然他们就只能掩耳盗铃地追求高正确率了，通常最终的结果就是“零取整存”。事实上，我曾经也经常受到隔夜跳空的伤害（一夜亏损6位数），而且我深信以后也无法避免受到这种伤害。在我看来，隔夜跳空只不过是一种噪声，使价格走势并非呈现连续波动的特征。那么理论上来讲，止损幅度就应该有一个下限，过小的止损幅度就是没有意义的。大的止损幅度自然就决定了仓位不能太重，这样一来我就用轻仓的策略来包容这种隔夜风险。

再次，是“怕赌情结”。一般来说，我同别人聊交易都是挺受人欢迎的。但是几乎所有的人都反感我将交易同赌博相提并论。如果市场上、甚至是现实世界中根本就不存在确定性的话，那么我们该如何应对呢？两种选择，要不我们就自己来制造确定性的幻觉，从而生活在虚假的信念中；要不我们就正视不确定性、承认不确定性、接受不确定性、研究不确定性、合理地控制不确定性。遗憾的是，更多交易者宁愿活在幻觉中，也不愿意面对不确定性。他们宁愿将自己的交易技术弄得模糊不清。这时交易者就不能使用确定的一致性交易方法，毕竟任何确定下来的交易方法都会立刻被证实不具有确定性。所以交易者只好毫无章法、随心所欲地完成交易行为。交易者天真地以为，众多具有不确定性的交易方法混合、搅拌在一起就能具有确定性。在我看来，我是一位知道自己是在赌博的交易者，而更多的人是自以为在做交易的赌徒！

最后，是“怕错情结”。如果谁能不犯任何错误，那么他将会快速地赚取整个世界。犯错并不可耻，逃避错误才是最可耻的！所谓“怕错情结”是指，不愿意主动承担责任的倾向。很多交易者在还没有行动之前，就找好了推卸错误的借口。其实明明已经做好了不可能改变的决定，但还要征求别人的意见。如果得不到他人的支持，那么就会继续征求其他人的意见，直到有人能够支持他的观点为止。这样的好处在于，当结果真的证实犯了错时，即使不能少亏一分钱，也能找到一个可以指责的对象。似乎对于他来说，自尊心比钞票更重要！

\subsection{圆木桶理论}

{\kaishu
对于一笔获利的交易而言，你必须把各个环节都做正确；但对于一笔亏损的交易而言，并不需要你把各个环节都做错，只要做错一个环节就足够了！

\hfill ——金融帝国
}

两杯50度的水，倒在一起不会成为100度的水。这说明，我们根本无法将温度叠加。如果我们承认，任何的交易方法都无法指向确定的结果，而只不过是基于概率下的产物。那么我们是否也同样无法叠加概率优势呢？

与此有异曲同工之妙的是“圆木桶理论”，木桶能够盛多少水，不取决于最长的那块木头，而取决于最短的那块木头。而你的木桶越复杂，就越容易出现更短的木头。事实上，交易也是完全一样的。一笔成功的交易需要正确的开仓、正确的持仓、正确的离市、正确的资金管理。而这其中一个环节出现问题，就可能出现交易最终的失败。没错，对于一笔获利的交易而言，你必须把各个环节都做正确；但对于一笔亏损的交易而言，并不需要你把各个环节都做错，只要做错一个环节就足够了！

我见过很多在期货市场明明看对行情，但由于没有做好资金管理反而输钱的交易案例。表面上看，期货既可以做空又可以做多，还可以使用资金杠杆，似乎获利会更加容易。但事实上，期货市场更容易让交易者死得更快、更惨！道理很简单，任何一个环节的问题，都可能让你被市场扫地出门。

如果说一笔成功的交易，需要众多的环节组成。而你的每一个环节又都需要（依赖）多种的交易技术来完成，那么使你交易行为出现问题的可能就会更多，从而最终失败的概率就会更大。一点都没错，交易的成败不取决于你最大的优势有多大，而取决于你最大的劣势有多大。这也是我反复提倡交易应该简单而不是繁琐的原因，毕竟交易是一种“失败者的游戏”。

以前做手机销售时，我发现手机的质量同功能成反比。或者说，功能越多的手机，就越容易出现故障。那么交易是否也会出现同样的问题呢？如果我们交易的系统大厦过于庞大与复杂，那么是否其中一个环节出现问题，整个交易的“流水线”就会瘫痪呢？

\subsection{点球原理}

{\kaishu
你需要知道的事情越少，获利就会越容易一些。

\hfill ——金融帝国
}

虽然我对足球并不算了解，但我老爸很喜欢看足球比赛。记得一次世界杯期间，我也受到流行潮流的影响而看起了足球。当一场比赛打成平局时，最后需要用残酷的点球大战来分出胜负。如此重大的赛事，我能够感受到选手那种巨大的心理压力。

我老爸对我说，射手的心理压力要比守门员大得多。我不理解其中的缘由，明明是一对一的对决，为什么会有一方的压力更大呢？老爸解释道：射手能否射好球，很大程度上取决于技术，所以担心发挥不好就有压力。而守门员却更大程度上取决于运气！也就是说，即使不能扑到球，也不容易遭到观众的责怪。

如果守门员等到射手已经将球踢出，甚至是看出球运行的线路再做反应的话，那么基本上就不可能扑到球。守门员必须在射手起脚的一刹那，赌球射向的方向，否则就一切都来不及了。事实上，我们经常能够看到点球大战中，球射向左面而守门员却扑向右面的场面。明白我说的意思吗？守门员扑点球不是看出来的，而是赌出来的。如果是赌左、中、右的话，那么扑到的概率为33％；但如果等看出球的走向后再选择扑球方向，那么扑到球的概率为0。

不追求确定性还有可能成功，而追求确定性就必然会失败！如果将守门员扑球比做交易的话，那么三分之一的正确率在五倍的赔率面前照样可以获利，而追求确定性的交易者通常什么都得不到！

曾经的我在捕捉趋势时，还需要其他方法的确认、需要自己一定能够看懂行情。后来我发现，也许迟钝（需要更多的条件相当于一种条件的迟钝）能够提高正确率，但它也是靠丧失赔率来换取的。而真正意义上的确定性又是不存在的，再有把握的交易也是可能出现亏损的。在我们追求有把握的同时，又将自己的成本大幅提高，从而丧失了赔率。

后来我发现，看出趋势的思维方式是错的！趋势不是看出来的，而是赌出来的。我现在的思维方式是：大众用潜在的好位置换取钞票，而我用钞票来换取潜在的好位置。请注意，我是用真实的钞票来换取潜在的好位置，所以我通常都在输。我从来都不买彩票，这倒不是因为彩票的成功率低，而是因为彩票具有负期望收益。即使彩票还是同样低的正确率，但如果我们能买到0.2元一注的彩票，那么投资彩票也是一条不错的生财之路。

相信我，我们从来都不可能以“零成本”知道更多的真相。所以说，你需要知道的事情越少，获利就会越容易一些。

\subsection{小蝴蝶的故事}

{\kaishu
少则得，多则惑。

\hfill ——老子
}

这是我在幼儿园时听过的一个故事。也许我们这些成年人，应该回过头来再听一听小时候的童话，或许能得到一种更深刻的感悟。

一只小蝴蝶觉得自己身上的色彩不如同伴的好看，就请求上帝帮自己改变身上的颜色。上帝满足了它的要求，给它一种自由选择颜色的权利。上帝指着五彩缤纷的花丛，对小蝴蝶说：“你喜欢什么颜色，就在什么颜色的花朵上‘打一个滚’，这时你就能够变成什么颜色。”小蝴蝶欣喜若狂地先是在红色的花朵上“打滚”，然后又不想放弃蓝色的花朵。这样就又在蓝色的花朵上“打滚”，这时它又觉得黄色的花朵美丽……就这样，它在所有的花朵中都“打了一个滚”，最终它变成了一只黑颜色的蝴蝶。它顿时痛哭流涕，无法接受自己成为一只颜色最丑陋的蝴蝶。

人们的痛苦往往不是来自于没有选择，而是来自于选择太多而无所适从。所谓选择就是一种利弊的取舍，这个世界上没有全优的选择，否则就谈不上什么选择了。当交易者面对众多充满特色的交易方法时，通常无法进行适合自己的取舍。似乎每一种方法都有让自己获利的时候，同样每一种方法也都有让自己亏损的时候。这时交易者可能就不做任何选择，而希望能够得到全部方法的优势。

你可以选择“赤橙黄绿青蓝紫”中的任何一种颜色作为自己的交易方法。但如果你吸取了所有的颜色，结果只能是什么颜色都没有的黑色！在叠加交易方法的过程中，不但是将全部方法的优势进行叠加，同时也会将全部方法的劣势进行叠加。

我很清楚大多数交易者的动机，他们只不过是逃避不确定性的存在。所有的优势等于可能捕捉到所有的利润，但所有的劣势又等于可能遭受所有的亏损（承担所有的成本）。事实上，如果一种交易方法能够具有明显的期望收益优势，那么虽然单次（微观上）无法找到确定性，但却能够实现长期战绩（宏观上）的确定性。有意思的是，交易者为了逃避微观不确定性，反而将自己带入了宏观不确定性的深渊中。

\subsection{损耗原理}

{\kaishu
趋势跟踪系统产生的大赔率在嘲笑众多小亏损的同时，也在嘲笑其他的分析方法以及我的思考与智慧。

\hfill ——金融帝国
}

在我们生活的真实世界中，存在着众多的极限。比如，速度的极限是（光速）每秒30万公里；绝对零度为$-$273.15摄氏度，等等。同样在我们虚拟的市场上，也存在着很多极限。比如，RSI指标的数值不可能超过100，交易正确率不可能超过100％，等等。

稍微懂点物理学的人都能知晓，极限是不可能达到的。并且越接近极限，每提高微弱的一步都需要更大的努力。回到我们虚拟的市场中来，同样也会发生这种越接近极限就会越困难的现象。比如，KD超过（小于）50后，就会出现涨跌对指标影响的不对称性。而且这种不对称性随着指标越接近100（0）就表现得越明显。如果KD指标达到95以上，那么我可以预期其后指标数值下跌的可能性更大。因为这时涨跌对于指标的影响已经非常不对称了，微弱的下跌就能使指标回落，而巨大的上涨也很难使指标提高一点点。这就是技术分析中所谓“背离”产生的原因。因为从指标的设计公式中，就存在一种向中轴（50）回归的吸引力。

如果我们把这种原理运用到交易方法的组合中，同样会出现这种损耗效应。根据“树形原理”我们可以得知，任何分析方法都是基于概率下的产物。而概率永远不可能达到100％，那么这时两种分析方法的优势在叠加过程中就会出现损耗现象。

概率越接近100％，提高一点就越困难，而降低一点就越容易。同时由于分析方法都是概率下的产物，那么一种方法的概率优势就会在另一种原有方法极大概率优势的面前出现损耗。或者说一种方法的概率劣势在另一种原有方法极大概率优势面前加倍放大。

趋势跟踪系统产生的大赔率在嘲笑众多小亏损的同时，也在嘲笑其他的分析方法以及我的思考与智慧。毕竟趋势跟踪系统的正确率通常在40％以下，那么靠其他技术来规避假信号将会变得非常容易。甚至我们随机的放弃一次信号，结果都有可能是避免了一次亏损。问题是，万一我们放过了一次真信号指出的大行情时。不但我们之前所做的全部努力带来的“好处”都会荡然无存，而且会出现得不偿失的结局。当趋势跟踪系统的赔率极大时，架构在趋势跟踪系统上的行情研判，就必须具备极高的正确率才有意义。想一想，如果你研判行情能够有极高的正确率的话，那么你还要趋势跟踪系统干什么？悲观地说，无论你使用趋势跟踪系统，还是自我的主观研判都能够实现不错的交易利润。但将他们组合起来的结果，很可能是不如单一使用其中的任何一种方法。

\subsection{依赖现象}

{\kaishu
把自己的钱袋寄托在别人的身上是一件最为愚蠢的事情。

\hfill ——金融帝国

你可以推卸责任，也可以走向成熟；但你无法既推卸责任又走向成熟。

\hfill ——金融帝国
}

我不知道是因为中国人的天性，还是因为我们所受到的教育，总之中国人更喜欢依赖外在的事物。这一点同美国人的风格截然相反。一位父亲对他的孩子说：“你从桌子上跳下来，我会接着你。”孩子真的跳下来了，但父亲却没有接着他，结果孩子摔得很痛。孩子指责父亲说：“你为什么骗我？”父亲说：“这个世界上你能够相信和依赖的人只有你自己。”这绝对不是一个简单的故事，事实上这确实是美国人风格的真实体现。与此有异曲同工之妙的是，美国上个世纪初关于实施“百日新政”探讨中的一段对话。弗兰克林准备动用国家资金购买过剩商品进行销毁，并且给每位储户的账户中增加一定数量的美元。反对者指出：“这会动摇我们引以为豪的‘民族责任感’。”

美国是一个讲究自由的国家，自然也就必须讲究每一位个体的责任感。中国则不同，小的时候我们就受到这样的教育：相信别人，相信政府；政府会为每一位公民负责，公民可以完全依赖政府。这样一来中国人就丧失了最为宝贵的责任感，我们似乎根本不知道什么是为自己负责。请别误会，我并非觉得相信政府有什么不妥，只不过这不应该成为大家逃避责任感的借口。否则，命苦就去怨政府，点儿背就去赖社会，就绝对不是什么好习惯了。

我有一个惊人的发现，中国股民会把某某事物当做“救世主”，又会把同样的事物当做“出气筒”。请注意，是同一事物！更有意思的是，这种事物会不断地被交易者当做救世主、出气筒而形成一种循环。

交易亏损了，就是股评的错、国家的错、上市公司的错。似乎自己从来都没有错！在行情低迷的2005年，各大论坛上怨声载道。有骂证监会的、有骂政府的、有骂上市公司的、还有骂股市黑嘴的。说真的，我还真的没有听到骂自己的！我就奇怪了，哪个国家的政府能够让股市只涨不跌；哪个国家的股市能够没有业绩虚假；哪个国家的分析者能够次次正确？

新股民通常把股评奉为神灵，然后就是相信股评、依赖股评、指责股评；事情还没有结束，然后是继续相信股评、继续依赖股评、继续指责股评；事情仍旧没有结束，然后是仍旧相信股评、仍旧依赖股评、仍旧指责股评。股民不光是这么对待股评，而是对待任何事物都是这样。

更有意思的是，股民在咨询他人意见的时候，通常是抱着事后推卸责任的动机。没有人愿意真诚地征求别人的意见，他们通常想要的不过是认同与支持罢了。如果第一个征求意见的人没能认同他的话，那么他会征求第二个人的意见，直到得到支持为止。如果最终自己做对了的话，那么他会认为是自己的高明。如果最后做错了的话，那么他会认为是受到别人的误导。有一点我百思不得其解，为什么几乎所有人都愿意把自己放在受害者的位置上呢？难道这是一种自虐倾向吗？难道受害者的位置真的这么具有吸引力吗？真正的强者不需要同情、真正的强者不会去指责任何人！

现在的我从不指责任何人！我相信强者的逻辑应该是：任何能够伤害到你的事物，而你又对这种事物无能为力时，那么这种事物绝对是没有任何错误的！只有当你能够使伤害到你的事物付出代价时，那么对方才能算是犯了错误。也许你不是强者，我允许你指责。但请记住，你只能指责一次、并且在指责过后就再也不能相信它了。有一句关于感情的名言需要大家记住：没有人能够伤害到你，除非是你自己愿意受到伤害。骂过上市公司业绩虚假的交易者，你们以后就不要看业绩报告了。我骂过上市公司的业绩虚假，而我又无法避免这种虚假的情况，所以我以后再也不看业绩报告了。我受到的伤害是值得的，毕竟我往前走出了一小步。

\subsection{绿草坪效应}

{\kaishu
请不要羡慕任何人，交易获利的关键在于：守住自己的“道”。

\hfill ——金融帝国
}

一山望着一山高，说的就是这个道理。或者说，在草原上你总会觉得远处的草坪比脚下的草坪更绿。我在生活中也发现过这种现象。在一条繁荣的步行街上散步，我从各个方向都能看到正在消费的人们。我似乎感觉到所有人的生活都比我优越，似乎人们都在不停地进行消费。哎，人家那样才能叫做生活。

后来我发现这里有一个重大的误区。我不是天天出来消费，那么我看到的这些消费者是否也并非天天出来消费呢？我时刻都能看到不停消费的人群，但他们事实上并非都是同一个人。可能我一天中只有少数的消费行为，但也能融入这个连续消费的场景中。我能看到不停消费的人群，但不等于人群中的人都在不停地消费。

我曾经听到有人说：“张三在1800点就指出上证指数能够涨到4000点，他真的是名副其实的分析大师。”我说：“如果大盘现在跌破1000点，你会说李四的观点太神奇了，他是名副其实的分析大师。”期货市场上有这么一句话：吹牛大旗永不倒，吹牛之人轮流换。

股市任何时候都能找到过去一段时间观点正确的分析师，并且媒体也乐于挖掘这种事后合理的现象，从而证明市场是能够被看懂的。请注意，总有观点正确的分析大师，并不等于有某位分析大师的观点总是正确。

我想说的是，任何市场参与者都有高潮和低谷，也都有“占便宜”和“吃亏”的时候。大家不要羡慕任何人，交易获利的关键在于：守住自己的“道”！别人赚钱就眼红、别人获利就模仿，那么你就会陷入“学习效应”的怪圈中。任何人使用的交易方法都是基于概率的，根本就不存在能够看懂行情的人，否则市场就会失去其存在的根基。

\subsection{焦虑原理}

{\kaishu
人的一生总是在追逐确定性，认为只有得到确定性才能使自己感到安全。事实上，确定性是不可能追求得到的。而只有意识到确定性并不存在时，我们才能使自己的内心得到安宁。

\hfill ——摘自《短线交易大师：工具与策略》
}

对确定性的追求，从某种角度上来看，也可以理解为对完美的追求。这是由我们每一个人的本能所决定的。如果我们失去确定性的话，那么就会陷入一种可怕的恐惧当中。从理智上来讲，不确定性是居于“得”与“失”之间的一种状态，它优于“失去”，但劣于“得到”。有意思的是，从心理学的角度上讲，不确定性比失去更为可怕。我相信大家绝对不认同我的说法，毕竟这又不符合逻辑！下面我就讲述几个真实的故事吧。

如果你赶12点的火车，但路途中却遇到交通堵塞。对于你来说，能否赶上这班火车是不确定的。这时你会非常焦虑，甚至会焦急得坐立不安。有意思的是，假如你最终知道火车已经开走了，你反而会定下神来，只是品尝着后悔。心理学指出，人们往往在事情发生时反倒不是很在意，而在事情可能发生的阶段最为焦虑。

“二战”期间有份研究报告曾指出，战斗机的飞行员非常容易焦虑，因为每一次出动都可能被击落。有些飞行员甚至想选择一死，如此就不用担心会被杀或受伤，他们索性抱着必死的信念去撞山。很多逃亡的罪犯最后选择投案自首，这并非他们良心发现，而是因为他们实在受不了可能落网的心理恐惧。他们宁愿选择自首来彻底消除可能落网的心理恐惧。我见过一些恋人，在一方因某种原因思考是否应该分手而犹豫不决的时候，另一方因实在受不了可能分手的焦虑不安，而主动选择放弃这段感情。

做期货的朋友可能会有与我同样的感受。我为了提前知晓第二天可能出现的盈亏，而情不自禁地在晚上关注外盘走势。其实外盘与国内市场并非百分之百相关，有时还有可能走出截然相反的走势，但大多时候却是同步运行。当每次行情做反，最难受的时候并非是第二天确定亏损了几万元，而是在前一夜只是知道第二天可能亏损的时候。有的时候这种焦虑真的让人难以入睡，而通常无论第二天损失多么惨重，我都能很快地接受并且适应。可能出现的亏损，比确定的大幅亏损更加可怕。

不可思议吧，人通常真的是不符合逻辑的。如果一个交易者不能战胜自己的心灵，那么他的行为可能更注重确定性而不是利润。人的一生总是在追逐确定性，认为只有得到确定性才能使自己感到安全。事实上，确定性是不可能追求得到的。而只有意识到确定性并不存在时，我们才能使自己的内心得到安宁。焦虑原理最为主要的一个提示就是，建议交易者放弃对不真实的确定性的追求。

这种“焦虑原理”通常能够形成一种巨大的恐惧，从而使交易者可能无法挺过交易系统中付出成本的阶段。即使我们的交易系统是采用保守性原则来设计合理的资金管理策略，从而保证账户资金的长期安全性。但我们仍旧不知道成本具体需要付出到什么时候，或者说还需要亏损多少钱才能熬过“交易的冬天”。这时交易者有可能在焦虑中，放弃未来能够给自己带来巨大利润的交易系统。正是这种也许明天就能够开始下一次的辉煌，也可能还要继续亏损一段时间的不确定，使得交易者心力交瘁。我们可以很轻松地在历史走势中测试系统的有效性，但可惜我们无法测试出心理反应对“按系统操作”可能带来的影响。当我们作系统测试时，可以问自己是否能够承受30％的资金回撤幅度，但我们无法设想那一时点上可能出现的心理问题。

\subsection{赌场原理}

{\kaishu
对于赌场而言，赌博的结果具有微观不确定性，以及宏观确定性。

\hfill ——金融帝国
}

赌博是一种“上不了台面”的行为，也是一种被大多数人所不耻的行为。那么有没有靠赌博为生的行业呢？在美国的拉斯维加斯、中国的澳门，都存在众多的赌场。而赌场的经营者确实能够靠其为生，甚至还能实现丰厚的必然利润。难道赌场的经营者不是在赌博吗？其实从某种意义上来讲，彩票的发行者也是靠赌博实现必然盈利的。

对于赌场而言，它与每位赌客的赌博结果是不确定的。但当它与足够多的赌客进行赌博后，就能产生必然的利润。这种利润来自于赌博规则对于赌场略微的概率优势。也就是说，对于赌场而言，赌博的结果具有微观不确定性，以及宏观确定性。甚至，当赌场与足够多的人次进行赌博后，最后的收益都能够通过数学公式精确地计算出来。并且赌博的次数越多，计算得出的理论数字就与最终的实际结果越接近。

从人类的科技发展史来看，早期的经典物理学，把世界描绘成一座精确的大钟。未来是可以通过绝对的因果关系进行预测的。牛顿在尝试股票交易后，曾经说过：“我能够精确地计算出天体的运行，但却无法计算人性的狂热。”事实上，几百年前天体的运行就能够被人类精确地计算。后来爱因斯坦提出的相对论，虽然使得经典物理学大厦崩塌，但是仍旧没有放弃因果关系。最近科学发现都在反复证明：对于观察者而言，能够预测的线性因果关系并非常态，而不可预测的非线性关系才是世界的常态。虽然线性因果关系能够解释简单系统，但仍旧有无法解释的两种情况。一个是复杂系统。即使天体再庞大，能够产生作用的因素却极其有限，因此它仍旧属于简单系统。对于复杂系统，可以用混沌理论加以描述。经典物理学虽然可以计算出天体的运行，但却无法计算出小溪流过岩石产生的波纹。另一个是微观世界，在微观世界中，对于电子运行轨道的描述就必须使用“波函数”，从而无法给出确定的结果，这就是著名的量子理论。简单地说，就是对于电子出现的位置而言，单次结果具有不可知性；但多次结果可以用“波函数”这种几率波进行精确的描述。有意思的是，量子理论与赌场的盈利模式具有惊人的相似，那就是：微观不确定性，以及宏观确定性。

我不得不感叹世界大同。原来最高深的科学与最低俗的赌博具有同样的道理。那么我们交易的市场是否也会出现这种现象呢？答案是肯定的！市场上没有任何方法能够指向唯一确定的结果，否则就会产生市场怪物。那么所有的交易方法都是基于概率的产物！这样一来，交易的盈利模式与赌场的盈利模式就别无二致了。所不同的是，赌场靠规则上的概率优势实现长期的必然盈利；而交易者靠人性的不对称性建立起来的概率优势实现长期的必然盈利。

\subsection{超级掠食者理论}

{\kaishu
如果你不知道未来市场会如何运行，那么你只要知道市场如何发展会对大多数人最不利，就等于得到了预知未来的水晶球。

\hfill ——金融帝国
}

人类的很多行为，往往都会被打上遗传心理学的烙印。我们祖先的生存策略，毫无疑问会在不同程度上影响现代人的行为。很显然，人类本能上的任何偏向都会在市场上体现出来。那么，我们又能从中得到一些什么宝贵的启发呢？

在一次同好朋友的闲聊中，得知他所从事的行业，赚钱有可能会变得越来越困难。本来工作上的不确定性并不是什么大不了的事情，但加上每个月必须偿还房贷的问题，就让他显得压力不小。事实上，他每个月的收入应该可以达到房贷还款的七八倍。就算真的有一天他要转行，估计凭他的经验也不会赚得很少，看来他的确是有些过于敏感了。相反，他对我在交易上的压力显得有些很不屑。他总是觉得，就算我亏损得还剩下个零头，在金钱上也要比他更有安全感。

接下来我跟他讲：房贷的压力并不是你应该考虑的事情。在贷款买房的人当中，如果只有你还不起贷款，那么这就是你个人的问题了；但如果有超过一半的人还不起贷款，那么这就是社会的问题。你替社会操心干什么呢？至于工作的问题呢，其实你也没有担心的必要。我以前做手机行业的时候呢，发现很多老板以前都是做家电的。当音响不赚钱的时候，这个圈子里的人都开始做传呼机；传呼机没落了没有多长时间，手机就火爆起来了；现在手机行业也没落了，但那些小老板不也都活得很好吗？事实上，你只要跟着大多数人形成的潮流往前跑，别人干什么、你就干什么，与此同时保证自己不是最差的就好了。而我就不同了，我是在没有航海图的水域中，努力开辟一条属于自己的航线。我遇到的所有问题，都是我所特有的，我没有可以依赖的环境。而你则不同，上可以依赖国家，下可以依赖你所从事行业的圈子。当然，这个世界也很公平。我如果成功了，那么得到的收获也是你所不能比拟的。你很难发大财，是因为你在任何环节都受不了比大多数人差，所以在事业上你只能做加法。而我的理念则不同，我相信穷人的努力不值钱，值钱的是穷人的牺牲。我为了将来能够享受更好的生活，所以可以心甘情愿地面对去过更差生活的可能性。问题是，你相对于我的优势不就是可以无忧无虑地轻松生活吗？如果你再庸人自扰，那么你的谋生手段的价值又是什么呢？

这就是“羊群思维”真实的写照。一只离群的小羊遇到可怕的掠食者时，通常很难有生还的可能。而当羊生活在羊群中时，只要保证自己不是跑得最慢的，就可以安全地躲过危机。这是一种有意思的生存策略，掠食者可以更容易地捕捉到食物，而猎物却可以靠牺牲最弱者来实现生存的目的。但这种生存策略也有一个致命的弱点，那就是可以使群体变得极端没有效率。毕竟，个体比拼的不是正确的方向，而是更快的速度！生活在群体中的个体，往往会丧失思考的能力。而集体形成的盲目性，又很容易将其带入毁灭的深渊。2008年美国的金融危机，就是羊群效应的真实体现。美国大量金融机构所犯下的错误，如果不是集体错误，那么毫无例外地都会经受灭顶之灾。对于基金经理来说，也时时刻刻都受到羊群思维的影响。很显然，即使是同样级别的错误，“特有的错误”与“集体的错误”绝对会受到投资者截然不同的惩罚（赎回资金）。而趋势跟踪交易者，从来不会追求相对（指数的）收益。追求绝对收益的另一面往往是，“错误”一定是特有的！恐怕，这就是趋势跟踪策略更难吸纳资金的原因。

与羊群类似的还有鱼群，但鱼群在超级掠食者面前就没有这么幸运了。人类的智慧可以建造出巨型的捕鱼船，甚至还可以通过超声波精确地定位出鱼群的位置。事实上，生活在鱼群中的鱼类已经开始变得越来越不安全。相反，离群的孤独之鱼，反而不容易被超级掠食者“惦记”上。

交易做久了，似乎让自己变得越来越像那只离群之鱼，从而使得自己对大众价值观有一种本能的排斥。遗憾的是，这已经成为我一种可怕的负担，这方面的话题我会在后面讨论。这里我想说的是：对于交易者来说，一定要相信市场上有超级掠食者的存在。我不想过多地讨论这种现象的成因，但效果甚至可以完美到成为预知未来的“水晶球”。

如果你不知道未来市场会如何运行，那么你只要知道市场如何发展会对大多数人最不利，就等于得到了预知未来的“水晶球”。在本书的开头，我用上一代人经历的经济现象来对比未来大尺度的经济走势，猜想未来房价如果涨到2000万该怎么办？事实上，按照这种理论来看，这种情况应该不会发生。对于上一代人来说，毫无例外的都是货币的储备者，而对于储蓄者最不利的情况就是通货膨胀。而我们这代人已经没有了艰苦朴素的储蓄习惯，甚至为了解决住房问题都背负了不少贷款。信贷扩张使得大多数人的角色已经从债权人转变为债务人！我相信，过去几年的信贷扩张、通货膨胀很可能在大尺度上使我们这代人成为“紧缩的一代”。大多数人手中的债务在增加、上一代人的存款已经被逐步消费掉，那么通胀的结束应该就不会太过遥远了。既然我们这代人手中没有什么存款、反而还背上了房贷，那么通货膨胀就会对我们这代人有利。我相信，对大多数人有利的事情，是一定不会发生的！

\subsection{卖完理论}

{\kaishu
对于供给面市场化严重不足的市场来说，供给的匮乏往往不会使价格上涨，反而会使价格崩溃。

\hfill ——金融帝国
}

前面已经说了很多所谓的理论、效应，现在快到这个章节结尾的时候了。本来我觉得这个所谓的“卖完理论”根本没有必要单独写成一节，毕竟这个结论并不严谨、也没有太多的实用性。但我始终觉得，正确的交易应该就是一种正确的思考。如果没有养成优秀的思维方式，那么恐怕没有人能够从技术、策略、系统的层面上做好交易。有的时候，我真的感叹：是什么可怕的力量使得人们变如此“无私”，又是什么恐怖的力量使人们轻易就丢掉了最基本的常识？

我们可以举一个最简单的例子，为什么这么多的中国人反对人民币升值呢？其中有一个最可笑的理由，那就是会使外汇储备缩水。问题是，人民币升值真的会使外储缩水吗？假设1美元可以在美国买到1公斤的黄豆，那么无论人民币的汇率是多少，外储都可以拿出1美元在美国购买1公斤的黄豆。真正使外储“对外”缩水的只有美国发生通货膨胀这一种原因，而人民币升值并不会使外储“对外”缩水。事实上，中国民间是（卖出商品后得到的）美元的卖方、人民币的买方；而央行则是民间的对手盘。请注意，美国只是单纯地支出美元来购买中国制造的商品，而美国人手中从来就没有人民币。事实上，人民币的升值只不过使外储“对内”缩水罢了。那么，我们又该如何理解这种“对内”缩水呢？最直接的理解就是，央行无法用外储从民间全部买回曾经抛售出去的人民币，因为人民币变得昂贵了。当然，我们也可以从另一个角度来理解。这种所谓的“对内”缩水，类似于个税起征点提高使财政收入“对内”缩水。税率起征点的高低，直接影响资源在社会与政府之间的分配。而人民币汇率似乎也能起到类似的效果。简单说就是，人民币的汇率越高，民间就会消费掉越多的外储，从而使政府可以动用的外储变得更少。反之，民间就更难消费掉外汇储备。我真是觉得不可思议：是什么力量使买入者不希望自己买入的东西价格上涨呢？这就好像，一位买入股票的交易者，指责股票过快上涨一样可笑！

又有人说了，人民币升值会影响中国的出口、拖累中国的经济。事实上，这种理由同样经不起推敲。前面的理由主要说的是花钱的问题。而这种理由主要是讲赚钱的问题。人民币升值会影响出口，这是无可争辩的事实，而让我不明白的是：为什么经济发展一定要依赖出口呢？一个显而易见的事实：财富的创造来自于生产而不是消费。如果出口停滞的话，中国的经济又会如何呢？我告诉你，中国马上就会陷入美国1929年时的困境。美国当时选择的是，印钞票购买过剩商品，然后将过剩商品填海（销毁掉）。中国的做法很类似，印钞票购买过剩商品，然后将过剩商品出口到美国，并将美元储备起来。请注意，无论是销毁还是出口，（在没有消费美元储备之前）对国内来说都是完全等效的。唯一不同的地方在于，美国的政策是一次性的，而中国则是持续多年、并且希望继续持续的。两者的区别体现在，美国人明白了穷人的贫穷会影响富人的致富，所以必须让穷人富有到可以消费过剩商品。而对于中国而言，贫穷不但不会阻碍富人的致富，甚至廉价劳动力还可以使中国制造更有竞争力。明白我的意思吗？一个是国民的贫穷是富人的负债；另一个却是国民的贫穷是富人的资本。这就是美国的可恶之处。是美国大量地进口中国制造的商品，使中国产生巨大的贫富差距、资产泡沫。中国人勤劳的生产行为只是得到了货币符号，而货币符号的意义就是更贵的房价。这与学生们更加勤奋的学习，只不过是在推高大学录取分数线，而绝对无法改变大学录取率非常类似。市场价格是资源分配的有效工具，而更高的人民币汇率就会迫使中国的制造业更多地考虑，如何“为人民服务”；而不是如何“为美国人服务”。如果你是制造业的老板，那么反对人民币升值我是可以理解的。但如果你只是一位穷人，那么我真不明白你的脑子中想的是什么了？

澄清一下，我并无意从战略的高度来讨论汇率问题的是与非，更没有主张人民币应该升值。毕竟，任何经济政策都会存在硬币的两面。而以上观点只不过是站在一位“自私的小老百姓”的视角来看待汇率问题。这就好像税收对纳税人永远都是痛苦的，但税收又有其明显的合理性。我所关心的角度，并非是税率多高最好，而是大众的思维方式。究竟是什么样的力量使纳税人反对税率的降低呢？或许结论是显而易见的：大众对很多最简单与基本的问题，却存在完全无效率的观点。另一个非常有趣的案例是，股民对于扩容速度过快的指责。

郎咸平曾经这样评价房地产市场：供给面市场化严重不足，需求面市场化严重过度。事实上，这个经典的评论，几乎适用于国内大部分的市场。最为极端的案例就是权证市场，供给方被严格地限定在需要解决全流通问题的上市公司和券商身上，而权证市场的参与者却是没有任何进入门槛的。从某种意义上来讲，中国特色的权证市场的使命，就是为了解决中国股市的全流通问题。而实现证券市场的全流通，又是一件具有历史意义的重大利好。我仍旧无意从宏观的角度上讨论权证市场的问题，而只是继续探讨大众无效率的思维方式，记得，管理层为了抑制权证市场过度投机，而推出“创设”权证的规则时，大部分中小投资者都指责这种增加供给的“创设”制度，而几乎没有人去质疑为何只有券商才有创设的权利？与此类似的是，在过去很多年的时间内，只要新股发行的速度加快，就会有大量的中小投资者指责这种圈钱行为。

这本来是一个显而易见的道理：对于需求方而言，供给方更加充分的竞争会使其受益。在国际市场上，任何市场的参与者都有权在缴纳足够的保证金后做空权证市场（创设权证）；而任何符合要求的企业都有权在证券市场上发行股票。事实上，很多中国企业都跑去纳斯达克上市，而美国股民并没有抱怨中国企业“圈钱”。能够阻碍企业发行股票的原因只有一种，那就是企业觉得当时的市场出价太过便宜。而对于中国的企业来说，上市甚至一度成为一种“福利”。很显然，对于供给的限制，会使得市场价格虚高。有意思的是，作为市场需求方的中小投资者，却非常拥护这种虚高的价格。请注意，对买方有利的绝非是高价格，而是能够上涨的价格！股票在1元到2元之间波动与在10元到20元之间波动，对于交易者是完全等效的，唯一不同的是股利增加十倍！或许这就是股市投资者不能全部分享经济增长果实的原因，毕竟虚高的股价严重地透支了未来上涨的空间。值得庆幸的是，到目前为止股票的发行已经开始越来越市场化了。

在讨论了大众无效率的观点后，接下来我们该进入正题了。所谓卖完理论是指：对于供给面市场化严重不足的市场来说，供给的匮乏往往不会使得价格上涨，反而会使得价格崩溃。我们不妨假设一种情况：如果北京城市建设的密度达到日本东京的水平，种种原因的限制使北京不再有新增的土地供给。那么，北京的房价应该是上涨呢，还是下跌呢？按照传统的思维方式来看，房价当然会上涨！因为土地已经变成了不可再生的资源，只有存量的楼盘可以在市场上转手。那么，供求关系就会使北京房价反复上涨。而我的结论，完全与传统的观点背道而驰。没错，我坚信房价会下跌！

我们不妨看一看2007年的股票市场，自从5·30调整以后大盘股与小盘股的走势出现了明显的分化。我们不难发现，在这段大盘股飙升的行情阶段中，有大量的央企巨无霸上市发行。这是中国规模最大的一批企业！有意思的是，这些企业上市一个就少一个。我们不难发现，大部分央企都已经完成上市，而此后央企大盘股的供给会趋于枯竭。当市场的大盘股没有新增供应时，大盘股的价格反而出现了崩溃！或许，有人会说整个中国股市都已经崩溃了。事实上，问题并非如此，大盘股的跌幅远远大于指数的跌幅，甚至小盘股指数在此后的2010年还创出了历史新高。我们想知道的是，这种现象是一种特例呢？还是一种普遍存在的规律呢？

老股民应该都有这样的市场共识：当前市场主要发行什么样的股票，什么样的股票就会成为市场中的热点。有意思的是，大部分股民都习惯性地将市场的下跌归罪为扩容速度过快。那么扩容过快究竟是否真的是市场下跌的罪魁祸首呢？记得在我刚进入股票市场的时候，曾经发现过一个“奇怪”的现象：B股的价格往往只有A股价格的几分之一。要知道，B股市场已经丧失了融资的市场职能。相当长一段时间以来，B股市场都没有出现一只新股发行。按理来说，没有扩容压力的B股价格应该高于A股才对呀，为什么反而是死水一潭呢？记得那时B股市场不但是价格低，甚至每天只会有几笔成交。在我的印象中，当时有位股评家指出：如果中国股市真的丧失了融资的职能，那并不是什么好事。而是会变得像B股市场一样死水一潭。

记得我以前说过这样一种观点：任何商品的普及过程都是一个价格不断下跌的过程，初期只有花费高价才能提前享受到这种新鲜的商品，而能够延期享受的人只需要付出更低的价格。遗憾的是，当年我把这种规律套用在房地产市场上。毫无疑问，不但具有商品属性，而且更多体现为资产属性的房地产市场并不适用这条法则。开发商的天职就是不断地销售商品房，而资产属性的商品并不适用于供需曲线的理论！没错，我的意思是，更高的价格有可能会形成更多的需求。确切地说，不断走高的价格，才会激发出更大的需求。

事实上，开发商能够实现销售最大化的方式只可能是让房价不断地小幅上涨。比如，开发商让房价以一个相对偏低的价格开盘，然后每隔几个月（甚至是一个月）都将售价调高百分之几（甚至更高）。请注意，如果这种模式都无法销售出更多的住房，那么即使让房价打折也无法实现更多的销售。毕竟，房价已经大幅上涨了好多年，脱离了房价继续上涨的假设，买房甚至不如租房。那么，开放商讲述“价格”的故事基本上是徒劳了，而只能够给大众讲述“趋势”的故事。从某种意义上来讲，开发商就成为了某个楼盘的庄家，不但吸引了有真实需求的购房者，同时也吸引了不少的投机者。只要开发商想继续销售新建住房，那么支持房价就是他们义不容辞的责任。房地产市场交易成本偏高、变现性极差，所以投机客更像是投资客，因此即使开发商使房价小幅上涨也不会承受过大的获利抛压。我还听说过一个更有意思的故事：一位房地产中介的人敲我亲戚家的门，问是否有卖房的打算？我的亲戚当然回答不卖呀。这时经纪人给出了一个至少比市价高出50％以上的价格。并且说如果有意出售的话，他们可以帮助寻找买家，并且在房产交易的网站上刊登出来……开始我始终不明白，这样的把戏到底动机何在？后来我渐渐明白事情的真相了：房产中介并没有打算从我亲戚的手中赚到任何利润，而只是合理合法地将那套住房高价地刊登出来。我相信，房地产中介一定有附近的房屋急需出售，所以才在网络上制造出几个很高的价格作为参照物。发现没有？真正的卖家有义务想尽办法来支撑价格！！

我所在的城市，有一个面积不大的老城区。虽然这个地块具有优越的地理位置，但十年前却是一个典型的贫民窟。后来通过平房改造陆续拆迁，大量建造商品房出售。现在来看，这个地块基本上成为天津最高档的居民区。我希望能够用这个案例来验证我的“卖完理论”，看看未来是否会符合我的预测？开发商在过去很多年间，不断地在这个地块盖楼销售。到目前为止，这个地块可盖楼的土地资源已经趋于枯竭。我想知道的是，此后该地块楼盘的价格是否会出现相对整体房价水平的走弱？我认为没有新增商品房的销售，就会使得无人继续支撑房价。而房地产开发商们，就会继续寻找新的地块、营造新的销售概念。当然，更多的人坚信最后的住房被销售一空后，地块的不可复制性会使房价不断走高。

与此类似的还有权证市场。很显然，权证市场存在的意义就是解决股权分置。而在没有全部解决这个问题之前，权证市场必须是活跃的。否则的话，股票持有者是无论如何都不会接受得到“废纸”的补偿方式。而当股权分置的问题被顺利解决以后，权证市场就失去了存在的意义，从此权证价格开始向废纸的价值靠拢。

正如前面所说，我对卖完理论探讨的本意，在于分析大众的思维误区。事实上，大众的思维误区还有很多很多。比如，我曾经听到过这样一种说法：现在什么东西的价格都在上涨，房价绝对不可能下跌，最多只是涨得慢一些罢了。其实，我的逻辑正好截然相反：如果零售物价永远不会出现快速上涨，那么房价似乎就永远不会下跌。相反，零售物价的快速上涨无疑会迫使管理层采用紧缩的货币政策。而零售物价对紧缩货币政策并不敏感，但资产价格却对货币政策非常敏感。很显然，房价属于资产价格的范畴，而通胀却会压低资产价格。

在探讨市场的问题时，我发现了一个有意思的现象。学理工科的人，往往受不了我模糊的因果关系，以及粗犷的逻辑性。而学哲学的人，却能够接受我用“相关性”来代替“因果关系”的思维方式。比如，加息对于股票市场是明显的利空。但在大尺度上，利率往往与股指同向波动。当然，我们也可以对这种现象进行归因。利率往往与经济环境同向波动，而股指对于经济环境更加敏感。问题是，对于很多微观现象来说，我们很难清晰地进行归因。甚至，我还习惯于运用“因果倒置”的思维方式。记得有一次和朋友探讨股市节后的走势时，我很随意地说了一句：节后还有那么多新股发行计划，所以股指走势应该不会太过难看。新股发行的速度往往与行情的好坏具有正相关性，但用新股发行速度来预测行情就很难用因果关系来解释了。或许正因为如此，学理工科的人才受不了我如此不严谨的推理方式，而他们的行为必须建立在严谨、富有逻辑的推理之上。这让我想起了我闹过的一个笑话：我对着电脑教一位朋友炒期货的方法时，翻到了早籼稻的走势图。由于这个品种刚上市，而“籼”字又不常用。我突然发现，不知道这个字该读什么？为了消除这种尴尬，我风趣地说了一句：“字可以不认识，但钱还是要赚的。”我想说的是，对于交易获利而言，很多时候并不需要过于严谨，更不需要你无所不知。

很遗憾，交易者的思维方式，并不像交易策略那么容易传授。何况，剖析大众共有的思维误区，又会给我带来巨大的压力。更为重要的问题是，我的思维方式并不一定是非常严谨的。有的时候，只是一种感觉，而不像交易系统那样具有清晰、一致性的规则。甚至有时，我自己都无法完全说明白我的整个思维过程。但另一方面，正确的思考对于交易者的获利又是如此的重要。这里我最多只能建议大家养成思考的习惯，思考的结论正确与否反而并不是最重要的。毕竟，在思维层面上能够有权决定对错的只有你自己。即使是对于市场而言，一次获利的交易也并不能说明，你的分析就一定是正确的。多一些质疑精神、少一些思维禁区，或许是优秀交易者所必备的素质。相对论其实并不复杂，而爱因斯坦正是思考了大家不愿意思考的“公理”。相对于学习而言，我更看重思考的能力。我希望读者朋友能够从质疑的角度阅读本书，至少我是不会在经过思考与验证之前相信任何观点。人云亦云的人，绝对成不了出色的交易者。我并不希望更多的交易者学习我的方法，或许我的方法对你根本就是无效的。我希望看到的是，本书能够让读者产生一些共鸣，甚至得到那怕是一点点的启发。最后我送给大家一句罗马哲学家爱比克泰德的名言：如果将你的身体交给一个陌生人任意处置，你一定会感到愤慨。那么，当你将自己的精神交给一个偶遇者任意处置时，你难道不感到羞愧吗？

\subsection{微宏互换}

{\kaishu
我们在微观上追求什么，我们就会在宏观上失去什么！

\hfill ——金融帝国

一些读过本书测试版的朋友向我反馈，我的部分观点存在明显的矛盾之处。比如，贯穿本书的一个观点就是：不要奢望找到市场的逻辑，放下思考而愚钝地坚持系统。而上一章节中的“卖完理论”，我却极力提倡独立思考的重要性，甚至认为出色的交易就是出色的思考。这其中似乎有不可调和的矛盾！

起初我并没有足够重视读者的这种疑问，而是简单地回答：我只是想客观地阐述自己的感悟，并不想刻意地证明什么。我甚至认为，一部好的作品是应该留给读者足够的质疑与思考的空间。但后来我渐渐发现，这种“不可调和的矛盾”中蕴含着更深层次的哲理。

\hfill ——金融帝国

踏着英雄的足迹，在我们以为会发现邪恶的地方，找到了上帝。在以为杀害了他人的地方，我们杀害了自己。在以为走出去的地方，我们将回到自己存在的中心。在以为一切都完了的地方，我们却与整个世界在一起。

\hfill ——《神话的力量》
}

这是《通往金融王国的自由之路》开篇的一段文字，而这本书恰恰是我走进系统之路的启蒙著作。多年以后我似乎终于能够理解这句话的寓意，但最初给我的感受只不过是简单的“事与愿违”而已。

当我们用追求利润、逃避亏损之心来面对交易的时候，我们得到的并非是利润而恰恰就是亏损。我在想，为什么追求利润之心给我们带来的是可恶的亏损呢？恰恰是，当我在“追求”（兑现）亏损、“逃避”（挂起）利润的时候，反而能够不断获得可观的利润。我似乎发现了一种“微宏互换”的原则：我们在微观上追求什么，我们就会在宏观上失去什么！我突然发现，我更多的地方都是与大众截然相反的：我（在宏观上）用刚性的策略来实现（微观上）弹性的交易方向；而更多的人是用弹性的策略来实现刚性的交易方向。嗯，没错。更多的人在交易上根本就没有一致性的原则，却在单独的某笔交易上执着地坚持固定的方向。事实上，这样的事例太多太多了。很多人希望在市场上快速获得成功，从而使得股民们学习空仓要比学习止损还要困难。结果就是，欲速则不达！赢就是亏；僵硬就是灵活；快就是慢……这好像与中国古代的道家思想不谋而合，道家的哲学思想告诉我们崇拜自然、无为而治。

弹性的方向与刚性的规则，很类似于庄子所说，外化而内不化（交易化而规则不化）。我所关心的是，弹性与刚性、化与不化之间是否有某种深层次的联系呢？或者说，弹性的交易方向，是否必然要靠刚性的交易规则来实现呢？如果我们打开过去几十年美国黄豆的走势图，就可以明显地发现一个超大尺度的箱体。有意思的是，如果你使用中期尺度的趋势跟踪系统来交易（测试）这张黄豆走势图，你会发现，自己是在大尺度上实现完美的高抛低吸！那么，宏观上的“高抛低吸”，是否必然要依靠微观上的“追涨杀跌”呢？再或者说，能够在微观问题上放弃思考的前提，是否又必然是对宏观问题的深刻思考呢？既然这是一个哲学的问题，那么我们还是跳出市场来讨论吧。在我看来，讲几个生动的故事，要远远比引经据典更容易说明问题。其实，现实中事与愿违的案例还有很多很多。

每天清晨的早餐需要：果农种植草莓，然后再由工厂制成果酱；粮农种植小麦，然后再由面包房制成面包；奶农生产牛奶，然后再被奶站送到我的家门口。市场经济的魔力，能够使我可以用低廉的成本享受到一顿并不算丰盛的早餐。但即使是一顿简单的早餐，也需要多人分工合作。当然，我的工作也会给更多的人提供以低廉的成本就可以享受的服务。如果我们虚化掉“钱”这种东西，那么市场经济真是一种高效率的，“人人为我、我为人人”的奉献社会。有意思的是，这一切美妙的前提，也就是市场经济的基石却是“利己行为”。宏观上的高效互助，要靠微观上的自私自利来实现。

上个世纪美国黄石森林公园的一场大火，几乎将整个森林烧毁。期间美国政府动用了大量的人力物力来控制火灾，但结果还是烧光了整片森林。事后科学家分析火灾造成如此惨重后果的原因，得出的结论却是惊人的。造成黄石森林公园大火直接的原因并没有什么特别，这种原因不过就是存在数亿年的闪电和大风。为什么相同的原因会产生不同的结果呢？为什么同样的原因产生的大火，在以前只需要用直升飞机灭火就可以解决，而这次却根本无法控制呢？

树木的成长和森林大火能够形成一种自然的平衡，这种平衡维持了亿万年。他们是不可分割的两种力量，不可能单独存在。事实上，不同规模的森林大火每年都会发生，只不过这次演化成为不可收拾的局面。有意思的是，造成这次罕见大火的根源，恰恰是长期以来人类灭火的行为。树木的密度超过自然调节能够承受的级别，所以火势就无法控制了。科学家通过计算机设计出一个模型：在一个固定的面积上，以某种固定速度随机长出一棵树；同时按照某种固定的速度随机往这个面积上投掷一颗火柴。当火柴落在空地上的时候，不会引起火灾；而当火柴落在树木的位置上时，将会引起火灾并且烧掉能与其相连的其他树木。这样的游戏就可以无休止循环下去。因为火灾本身降低了其后火灾发生的概率，同时树木的成长又会增加火灾发生的概率。从某种意义上来讲，正是森林管理员长期的灭火行为，造成了这次毁灭性的火灾。此后，森林管理员有计划地、可控地烧毁过于茂盛的部分森林。

自然界有一种奇怪的特性，那就是遇强则强、遇弱则弱。当我们用一种对抗与征服的动机对待自然时，自然的力量会随着我们力量的增长而不断增长。抗生素会使细菌变得更强大，杀虫剂可以使害虫更强大。我曾经听过一种骇人听闻的学说：随着医学科技的发展，很多危及生命的疾病在现代医学的帮助下能够得到“治愈”。这就使很多本无法活到繁衍后代阶段的人能够繁衍他们的子孙。而在人类祖先的时代，他们都会在繁衍后代前死亡。现代医学在试图摆脱自然竞争、优胜劣汰的情况下，使人类种群基因变得越来越差。最终，没有经历过优胜劣汰的基因，有可能将人类整个种群变得更脆弱、甚至是毁灭。

这个话题似乎太过沉重了，沉重到我们无法探讨解决之道，至少我们在逻辑上无法质疑这种说法。在人类与自然的竞赛中，似乎我们已经没有停下来的可能。或许，有一天人类能够通过修改自身的基因来解决这个问题吧。咱们不要把话题扯远，还是继续回到交易当中来。人性中趋乐避苦的本性，无论是在火灾、疾病、还是亏损面前都表现得毫无二致。或许，我们在微观上的追求，的确会在宏观上造成某种隐患。

我始终觉得，出色的交易者多少要具备一些简单的心理学知识和哲学素养。对于这个话题的讨论，已经越来越接近道家的思想了。有兴趣的朋友，可以读一读《道德经》《庄子》。当你能够用辩证的角度来看待世界时，你一定会感叹世界的美妙。

当我在思考交易的问题时，我经常会运用“等效原则”“互损原则”。比如，我认为趋势跟踪系统中参数的敏感与迟钝是等效的；而“策略与技术”“正确率与赔率”，甚至是“心灵与效率”又是互损的。当你找到一种适合自己的思考方式后，市场将会变得简单得多。

记得有一次和朋友讨论学习（上学）与思考（在实践中自由感悟）哪个更有价值时，那位朋友说过这样一句话：我不认同你的结论，但我也要承认你的思维的确很美。或许，市场上从来就没有正确的答案，而有意义的只是优美的思考。因为答案正确与否，取决于观察者的角度与层次。接下来，我鼓起勇气与大家尝试探讨一下道家思想眼中的“善恶观”。我并不希望大家认同任何结论，而只是欣赏一下道家的思维方式。

看过《西游记》的朋友或许会发现，道观中的道士往往都是反面角色。这种现象恐怕来自于大家对于道家的误解，认为道家思想是一种自私自利的思想。我们该如何理解杨朱的那句“拔一毛而利天下，不为也”呢？其实，杨朱后面还有一句：“人人不损一毫，人人不利天下，天下治矣。”谁说市场经济是西方的舶来之物？这不就是对市场经济最早的主张吗！

或许，我们不难发现，道家的主张并非是“善良”、而是“公正”。那么，我们该如何定义所谓的“善”呢？最直白的理解就是：付出而不求回报。但似乎又有一些不太对劲的地方，毕竟很多人用“积德”之心去“行善”。很多人相信善有善报，就算这种福报无法降临在自己身上，也可以为自己的后代“积德”。或许，我们可以把关于“善”的定义再修改一下：付出了而暂时不求回报。

大家发现没有，抱着这种动机去行善，已经与市场经济下的储备行为无异了。当然，这需要假设你赚钱的方法是合法合理的，或者说是使用公正的方式来赚钱。如果自私之心是人的本性，那么在公正的前提下，“积德”与“存钱”就是完全等效的。

道家思想精妙之处就在于，他可以用辩证的视角来观察硬币的两面。善与恶、节约与浪费、生产与消费，必须是绝对相等的。这就像期货市场上，多头与空头的持仓数量必须是严格对等的一样。同时做多的行为可以推高价格，从而使继续做多变得更加困难，而让继续做空变得成本更低。如果每个人都奉献而不索取的话，那么这个世界会变成什么样子呢？很显然，那将会出现美国1929年那样的经济危机！每个人都为了“福报”而提供商品，但每个人都怕遭到“恶报”而无人问津。大家想想看，如果储蓄是一种善良的行为的话，那么行善的前提必须是有人“作恶”。难道不是吗？中国的过度储备离开了美国的过度消费，是不是根本难以为继呢？计划经济时市场中每个个体都将节约视为美德，但这却是一种经常发生资源错配、浪费严重的经济模式。市场经济正好相反，把消费视为拉动经济的法宝。市场中每个个体都在尽情“浪费”，但资源最佳配置却使整体经济表现出最为节约的效果。

按照这种逻辑，微观上行善的意义就是在宏观上纵容他人作恶，或者说让“恶人”有更容易作恶的土壤。前些日子看到腾讯新闻中关于“残疾儿童乞讨”的报道，其中讲述的内幕简直令人发指！难道不是因为善良人们的施舍行为，让这种可怕的恶行有了存在的土壤吗？道家思想精妙的地方就在于看到了：没有了善，恶就没有了土壤！

我保证，虽然上面的思维方式在逻辑上找不出明显的缺陷，但读者朋友们恐怕无论如何都无法接受。当然，我的说法也有一些断章取义的嫌疑。事实上，人性当中的确有善的一面，而这种付出是完全不求回报的。比如，儒家所说的恻隐之心。我家从小养了一条罗汉鱼，但这条罗汉鱼就是不起头。不起头的罗汉鱼，是没有任何经济价值的，甚至观赏价值也会大打折扣。目前还没有让罗汉鱼必然起头的技术，而做卖鱼生意的渔场就会将不起头的鱼杀死。那我该怎么办呢？扔了它再买一条？不！这样做的话，我心里会很难受的！毕竟，它是一条小生命呀！这就是恻隐之心，其实这种善良也很有意思：不是为了得到福报而行善，只是为了取悦自己的心理感受。如果将愉悦心灵也视为一种价值的话，那么人性仍旧是自私的。或许，这就是道家所说的，刻意去行善就不是真的善。否则，就会出现用“存钱”之心去“积德”的尴尬。

这应该是本书唯一一次质疑大众思维方式的地方。我说这些并没有别的意思，只是告诉大家：思想是不应该受到任何束缚的！同时，也让大家身临其境体会一下批判性思维能够带来的那种颠覆性震撼。仅此而已，毕竟这不是一本关于人生哲理的读物。

\clearpage
